% Source-specific reconstruction. ED26-SRC-027-501; not the 1992 chapter.
% Partial authoring block: cited results retain full-proof obligations.
\section{Klop's 1987 reduction-system tutorial}
\label{sec:p12-src-027-501}

This presentation follows J. W. Klop's CWI Note CS-N8701,
\emph{Term Rewriting Systems: a tutorial}, May 1987
\parencite{Klop1987CWITutorial}. Its objects are finite first-order terms
and abstract reduction relations. The report deliberately does not develop
conditional term rewriting or rewriting modulo equations. Its uses of
equations in completion do not change that boundary: completion seeks an
ordinary, directed rewriting presentation of an equational theory.
The six source sections have a definite dependency order: abstract reduction
properties; the term structure; a marked-tree termination order; completion;
regular systems; and normalizing or cofinal strategies for regular systems.
The 1987 report is a distinct presentation, not an earlier page layout of the
1992 \emph{Handbook of Logic in Computer Science} chapter.

\subsection{Abstract reduction: properties and their distinctions}

\begin{hbdefinition}{def:p12-klop1987-ars}{The report's abstract reduction vocabulary}
An abstract reduction system is
\(\mathcal A=\langle A,(\to_\alpha)_{\alpha\in I}\rangle\), a set with
an indexed family of binary relations. The report writes
\(\twoheadrightarrow_\alpha\) for the reflexive transitive closure,
\(\to_\alpha^+\) for the transitive closure, \(\to_\alpha^=\) for the
reflexive closure, and \(=_\alpha\) for the equivalence relation generated
by reduction. Thus \(=_\alpha\) is conversion, whereas an unadorned equality
between elements, outside a declared conversion formula, is identity.
Converse arrows reverse direction. Composition is read in traversal order:
\(a\to_\alpha\circ\to_\beta b\) means
\(a\to_\alpha c\to_\beta b\) for some \(c\).
The union of two relations is \(\to_{\alpha\beta}\).

Relations \(\alpha,\beta\) commute weakly when each peak
\(c\leftarrow_\beta a\to_\alpha b\) has a join
\(c\twoheadrightarrow_\alpha d\leftarrow\!\!\leftarrow_\beta b\).
They commute when their reflexive transitive closures commute weakly.
For a single relation, weak confluence or weak Church--Rosser (WCR)
means joinability of every one-step peak; confluence or Church--Rosser
(CR) means joinability of every finite-reduction peak.
Subcommutativity, \(\mathrm{WCR}^{\leq1}\), requires joins by at most
one step on each arm.
\end{hbdefinition}

\begin{hbdefinition}{def:p12-klop1987-properties}{Normal forms, bounds, and branching}
For a single-relation ARS, a normal form has no outgoing step, including
no self-loop. WN means that every element has a normal-form descendant;
SN means that no infinite reduction starts at any element. The converse
relation's termination is denoted \(\mathrm{SN}^{-1}\).
UN means that two normal forms reachable from the same element are identical.
NF means that a normal form convertible with an element is reachable from
that element. Ind means that every finite or infinite reduction sequence has
one element reachable from every member of the sequence.
Inc means that there is a map \(a\mapsto |a|\) into the natural numbers
strictly increasing at each reduction step; here the report's bars are a
natural-number ranking, not cardinality. FB means that every element has
finitely many one-step successors; \(\mathrm{FB}^{-1}\) applies this to
predecessors. A complete system is CR and SN. The report calls a system
with UN and WN semi-complete, while noting that this terminology was not
standard in its setting.
\end{hbdefinition}

\begin{hbdefinition}{def:p12-klop1987-subars}{Closed subsystems and reduction graphs}
A sub-ARS has an underlying subset, the restricted reduction relation, and
closure under successors in the larger system. An arbitrary induced
subgraph without successor closure is not a sub-ARS in this sense.
The reduction graph \(\mathcal G(a)\) consists of all descendants of
\(a\), including \(a\), with the restricted relation. A subset \(B\)
is cofinal in an ARS if every element reduces to some member of \(B\).
The cofinality property CP requires a finite or infinite reduction from
each \(a\) whose members are cofinal in \(\mathcal G(a)\).
\end{hbdefinition}

\begin{hbtheorem}{thm:p12-klop1987-ars-core}{proved=annex:p12-klop1987-ars-core}{Core implications in the original abstract route}
CR implies NF, which implies UN. UN together with WN implies CR and Ind.
Ind together with Inc implies SN. For a countable ARS, CR is equivalent
to CP. CR is also equivalent to joinability of every pair of convertible
elements, and to the mixed peak condition
\(c\leftarrow a\twoheadrightarrow b\Rightarrow
c\twoheadrightarrow d\leftarrow\!\!\leftarrow b\).
The report states these in Proposition 1.2, Theorem 1.6(i), (iii)--(v),
(vii), and Exercise 1.7(13), pp.3--8.
\end{hbtheorem}

Newman's implication SN and WCR imply CR is Theorem 1.6(ii); its
well-founded joining argument is supplied in
Theorem~\ref{thm:p12-newman} and its reciprocal annex. The additional
Theorem 1.6(vi), WCR and WN and Inc imply SN, is attributed in the report
to Klop's 1980 work. Unlike Newman's lemma, it cannot be obtained simply
by assuming SN: termination is its conclusion. Figure 2 also records
\(\mathrm{WCR}^{\leq1}\Rightarrow\mathrm{CR}\),
\(\mathrm{Inc}\Rightarrow\mathrm{SN}^{-1}\), and
\(\mathrm{FB}^{-1}\mathbin{\&}\mathrm{SN}^{-1}\Rightarrow\mathrm{Inc}\).
Its arrows have conjunctive endpoints; CR alone does not imply WN.
CR together with two distinct normal forms implies consistency, as
Exercise 1.7(14) explains.

\begin{hbdefinition}{def:p12-klop1987-ars-exercises}{The complete additional ARS problem surface}
Exercise 1.7, pp.6--8, uses ``exercise'' also for extensions of the theory.
The following preserves its seventeen separate claims and requests, not
only the ones used later in the report. Write \(\to_i^*\) in the
formulas below as a typographic expansion of the source's double arrow.
\begin{enumerate}
\item Rosen's first transfer: if \(\to_1^*=\to_2^*\) and
\(\to_1\) is subcommutative, then \(\to_2\) is confluent.
\item Hindley's union theorem: if each pair of indexed relations
commutes, their union is confluent. The case of a relation paired with
itself is included in the hypothesis.
\item Hindley's one-step condition: if
\(a\to_1 b,a\to_2 c\) can be joined by
\(b\to_2^*d,c\to_1^=d\), then the relations commute.
\item Staples' condition: if the same peak joins by
\(b\to_2d,c\to_1^*d\), then the relations commute.
\item Rosen's request relation: \(\to_1\) requests \(\to_2\) if
each peak \(a\to_1^*b,a\to_2^*c\) has
\(b\to_2^*d,c\to_1^*e\to_2^*d\). If each relation is confluent
and the first requests the second, their union is confluent.
\item The local mixed-peak condition
\(a\to_1b,a\to_2c\Rightarrow
\exists d,e\ b\to_2^*d,c\to_1^*e\to_2^*d\)
does not by itself imply full requesting, even when \(\to_2\)
is terminating and confluent. The counterexample is specified in
Theorem~\ref{thm:p12-klop1987-ars-countercases}.
\item Under confluence of \(\to_2\) and requesting, define
\(\to_3=\to_1^*\circ\to_2^*\). If every \(\to_1^*\)-peak
has both arms joining by \(\to_3\), the union is confluent
(Staples).
\item A refinement has \(\to_1\subseteq\to_2^*\). It is compatible
if each \(a\to_2^*b\) has a common \(\to_1^*\)-descendant of
\(a,b\). Prove that compatibility is equivalent to the property:
\(a\to_2b\to_1^*c\) implies some common \(\to_1^*\)-descendant
of \(a,c\).
\item Compatible refinement preserves and reflects confluence (Staples).
\item Strong confluence means that each one-step peak joins by
\(b\to^*d,c\to^=d\). Prove that it implies confluence (Huet).
\item Assume all indexed relations commute weakly. A relation is
relatively terminating if no infinite union reduction contains infinitely
many of its steps. It has splitting effect if a peak with that step on
one arm has no weak-commutation join using at most one step on the other
arm. Prove union confluence if every relation with splitting effect is
relatively terminating.
\item Le Chenadec's decreasing-conversion criterion: suppose a strict
well-founded order \(>\) contains every reduction step. CR is equivalent
to the requirement that every peak \(a\to b,a\to c\) admit a
conversion \(b=d_1\leftrightarrow\cdots\leftrightarrow d_n=c\)
with \(a>d_i\) for every \(i\), including both endpoints.
\item Cofinality and CP are as in
Definition~\ref{def:p12-klop1987-subars}; prove CR equivalent to CP
for countable ARSs (Klop). This proof is supplied in
Annex~\ref{annex:p12-klop1987-ars-core}.
\item Consistency means that not every pair is convertible. Show that
CR and two distinct normal forms imply consistency. A conservative
extension \(\mathcal A\subseteq\mathcal B\) preserves exactly the
conversion relation on \(A\). Show that a conservative extension of
a consistent ARS is consistent, and that a confluent closed extension is
conservative.
\item Newman's strict nontrivial diamond \(\mathrm{WCR}^{1}\) joins
each one-step peak with distinct endpoints by one genuine step on each
arm; equal endpoints already coincide. Prove
\(\mathrm{WCR}^{1}\mathbin{\&}\mathrm{WN}\Rightarrow\mathrm{SN}\),
and refute the version using only \(\mathrm{WCR}^{\leq1}\).
\item The quasi-commutation assumption is
\(a\to_\alpha b\to_\beta c\Rightarrow
\exists d\ a\to_\beta d\to_{\alpha\beta}^*c\), attributed to
Bachmair--Dershowitz. This assumption alone does not equate
termination of the two relations: on one
element take an \(\alpha\)-self-loop and empty \(\beta\).
The premise is vacuous, \(\beta\) is SN, and \(\alpha\) is not.
\item Let \(\iota:A\to B\), \(\kappa:B\to A\), with
\(\kappa\iota(a)=a\). Suppose every \(a\to_\alpha a'\)
lifts from every \(b\) with \(\kappa(b)=a\) to a genuine
\(b\to_\beta b'\) with \(\kappa(b')=a'\). Prove that SN
of \(\mathcal B\) implies SN of \(\mathcal A\) (Klop).
\end{enumerate}
Figures 4--5 distinguish genuine steps, reflexive steps, finite
reductions, conversions, and lifting maps; these distinctions are part
of the stated hypotheses. The following results supply the separate
proofs and countercases of these abstract assertions.
\end{hbdefinition}

\begin{hbtheorem}{thm:p12-klop1987-ars-commutation}{proved=annex:p12-klop1987-ars-commutation}{Finite strips, strong confluence, and commuting unions}
The conclusions of Exercises 1.7(1)--(4) and (10) hold with their
displayed arrow bounds. In particular subcommutativity implies CR,
equal reflexive transitive closures transport CR, and a family of
pairwise commuting relations has a confluent union, including when
the index set is infinite. The self-pairs in that family hypothesis
are essential: each individual relation must commute with itself.
\end{hbtheorem}

\begin{hbtheorem}{thm:p12-klop1987-ars-request}{proved=annex:p12-klop1987-ars-request}{The two full requesting criteria}
Under the full requesting condition of Exercise 1.7(5), the union is
CR if both relations are CR. More generally, under requesting and
CR of \(\to_2\), it suffices that every \(\to_1^*\)-peak
join by \(\to_3=\to_1^*\circ\to_2^*\) on both arms, as
in Exercise 1.7(7). This does not replace full requesting by the
insufficient one-step premise of Exercise 1.7(6).
\end{hbtheorem}

\begin{hbtheorem}{thm:p12-klop1987-ars-refinement}{proved=annex:p12-klop1987-ars-refinement}{Compatible refinement at the finite-reduction scope}
For a refinement, compatibility for \emph{every finite}
\(a\to_2^*b\) is equivalent to the mixed condition in Exercise
1.7(8), and it preserves and reflects confluence as in (9).
Compatibility only for single steps is insufficient for that equivalence.
\end{hbtheorem}

\begin{hbtheorem}{thm:p12-klop1987-ars-splitting}{proved=annex:p12-klop1987-ars-splitting}{Relatively terminating splitting colours}
Suppose every indexed relation commutes weakly with every other one,
and each relation with splitting effect is relatively terminating in
the whole union. Then the union is CR, as in Exercise 1.7(11).
An arbitrary finite peak uses only finitely many colours; relative
termination is used on that finite colour set, not silently converted
into termination of an infinite union of terminating relations.
\end{hbtheorem}

\begin{hbtheorem}{thm:p12-klop1987-ars-decreasing}{proved=annex:p12-klop1987-ars-decreasing}{Decreasing conversions give confluence}
The equivalence in Exercise 1.7(12) holds: for a strict well-founded
order containing all steps, CR is equivalent to a conversion joining
each one-step peak entirely below its source. The order must be
transitive, and the bound applies to every listed conversion element,
including its endpoints. An arbitrary conversion without that bound
does not supply this criterion.
\end{hbtheorem}

\begin{hbtheorem}{thm:p12-klop1987-ars-boundaries}{proved=annex:p12-klop1987-ars-boundaries}{Closed extensions, consistency, and genuine-step lifting}
CR, WCR, subcommutativity, WN, SN, UN, NF, Ind, Inc and FB all pass
to successor-closed sub-ARSs as defined above.
A confluent successor-closed extension is conservative; a conservative
extension of a consistent ARS is consistent; CR with two distinct normal
forms implies consistency, as in Exercise 1.7(14). The retraction and
genuine-step lifting hypotheses in (17) transport SN from the larger
system to the smaller. A zero-step lift would not prove that assertion.
\end{hbtheorem}

\begin{hbtheorem}{thm:p12-klop1987-ars-normalization}{proved=annex:p12-klop1987-ars-normalization}{The strict-diamond and increasing-rank termination arguments}
The strict nontrivial diamond together with WN implies SN, as in
Exercise 1.7(15). Also WCR, WN and Inc together imply SN (Theorem
1.6(vi)); Inc implies \(\mathrm{SN}^{-1}\); and
\(\mathrm{FB}^{-1}\) with \(\mathrm{SN}^{-1}\) implies Inc.
The strict-diamond result does not hold when each joining arm may
independently take zero steps.
\end{hbtheorem}

\begin{hbtheorem}{thm:p12-klop1987-ars-countercases}{proved=annex:p12-klop1987-ars-countercases}{The actual local-requesting and termination countercases}
The local premise of Exercise 1.7(6) does not establish full requesting;
the quasi-commutation premise in (16) does not equate SN of its two
relations; and \(\mathrm{WCR}^{\leq1}\) with WN does not imply
SN. Their explicit finite countermodels appear in the reciprocal annex.
These are nonimplication results, not assumptions of the positive theorems.
\end{hbtheorem}

\begin{hbexample}{def:p12-klop1987-abstract-failures}{Hindley's peak and van Glabbeek's puzzle}
Figure 3, p.6, has four elements \(a,b,c,d\) and exactly the steps
\(b\to a,b\to c,c\to b,c\to d\). It is WCR: the peak at
\(b\) joins at \(a\), and that at \(c\) joins at \(d\).
It is not CR, since \(b\twoheadrightarrow a\) and
\(b\twoheadrightarrow d\), with \(a,d\) distinct normal forms.
This is the counterexample attributed to R. Hindley.

The answer to Puzzle 1.6.1 on p.35, attributed to R. van Glabbeek,
can be specified without losing its infinite diagram. For each
\(n\in\mathbb N\), take \(u_n,v_n,z_n\), with steps
\(u_n\to u_{n+1},u_n\to v_n,v_n\to u_n,v_n\to z_n\)
and no others. Every element reaches a normal form. Peaks at \(u_n\)
join at \(u_{n+1}\); peaks at \(v_n\) join at \(z_n\).
But \(u_0\to u_1\to\cdots\) has no common descendant for all its
members: a fixed element belongs to a fixed index, and steps never decrease
that index. Thus WCR and WN do not imply Ind.
\end{hbexample}

\subsection{Finite terms, application, and disjoint signatures}

\begin{hbdefinition}{def:p12-klop1987-terms}{The source's ordinary first-order regime}
A signature \(\Sigma\) has a countably infinite variable set and a nonempty
set of operators, each of a fixed finite arity; arity zero gives constants.
\(\operatorname{Ter}(\Sigma)\) is the set of finite terms.
\(\operatorname{Ter}_0(\Sigma)\) contains the ground terms. A term is
linear when no variable occurs more than once. A one-hole context
\(C[\ ]\) gives the subterm relation \(t\subseteq C[t]\).
Subterm occurrences, rather than only subterm values, are distinguished
by positions: the root has position \(\epsilon\), and child indices
are concatenated along a path.

A substitution \(\sigma\) preserves operators and is determined by its
values on variables. The report uses \(t^\sigma\) for its action.
A rule \(l\to r\) has nonvariable \(l\) and
\(\operatorname{Var}(r)\subseteq\operatorname{Var}(l)\).
Its instances act in any context:
\(C[l^\sigma]\to C[r^\sigma]\).
The associated ARS has all terms as its universe. The notation
\((\Sigma,R)_0\) instead restricts that ARS to ground terms.
This differs from a ground TRS, in which every rule itself is ground.
The report keeps the single-sorted regime, noting the usual sorted
extension but not developing a new many-sorted calculus.
\end{hbdefinition}

\begin{hbexample}{def:p12-klop1987-arithmetic}{Source Table 1 and the ground/open distinction}
The signature has binary \(A,M\), unary \(S\), and constant \(0\).
Its four rules are
\[
 A(x,0)\to x,\qquad A(x,S(y))\to S(A(x,y)),\qquad
 M(x,0)\to0,\qquad M(x,S(y))\to A(M(x,y),x).
\]
For example, \(A(M(x,y),y)\) is nonlinear,
\(A(M(x,y),z)\) is linear, and
\(A(M(S(0),0),S(0))\) has two occurrences of the subterm \(S(0)\).
The p.10 calculation, writing \(2=S(S(0))\) only in this display, is
\begin{align*}
 M(2,2)&\to A(M(2,S(0)),2)
 \to S(A(M(2,S(0)),S(0)))\\
 &\to S(S(A(M(2,S(0)),0)))\to S(S(M(2,S(0))))\\
 &\to S(S(A(M(2,0),2)))\to S(S(A(0,2)))\\
 &\to S(S(S(A(0,S(0)))))\to S(S(S(S(A(0,0)))))
 \to S(S(S(S(0)))).
\end{align*}
Adding \(A(x,y)\to A(y,x)\) makes the open term \(A(x,y)\)
non-normalizing. Ground terms still have normal forms, but may now have
infinite reductions, including repeated swaps. Figure 6 records this
difference, rather than equating open and ground termination.
\end{hbexample}

\begin{hbexample}{def:p12-klop1987-sts}{The two string-to-term correspondences}
The string rule \(aba\to bab\) becomes
\(a(b(a(x)))\to b(a(b(x)))\), with unary \(a,b\) and constant \(o\).
The word \(bbabaaa\) encodes as
\(b(b(a(b(a(a(a(o)))))))\) and its reduct \(bbbabaa\) as
\(b(b(b(a(b(a(a(o)))))))\). Contexts encode prefixes and substitution
for \(x\) encodes suffixes, giving the source's isomorphism with the
ground-term reduction graph. Treating letters instead as constants and
using a binary associative concatenation operator would require rewriting
modulo associativity; the report explicitly does not take that second
route. Both sides of the unary rule retain \(x\), so these string-derived
TRSs are non-erasing.
\end{hbexample}

\begin{hbdefinition}{def:p12-klop1987-applicative}{The three displays of Combinatory Logic}
Application is the binary operator \(\operatorname{Ap}\); \(S,K,I\)
are constants. Source Tables 2--4 display the same three rules, respectively
in full binary, dotted, and left-associated notation:
\begin{align*}
 \operatorname{Ap}(\operatorname{Ap}(\operatorname{Ap}(S,x),y),z)
 &\to\operatorname{Ap}(\operatorname{Ap}(x,z),\operatorname{Ap}(y,z)),\\
 \operatorname{Ap}(\operatorname{Ap}(K,x),y)&\to x,\qquad
 \operatorname{Ap}(I,x)\to x;\\
 ((S.x).y).z&\to(x.z).(y.z),\quad(K.x).y\to x,\quad I.x\to x;\\
 Sxyz&\to xz(yz),\qquad Kxy\to x,\qquad Ix\to x.
\end{align*}
Juxtaposition associates to the left; parentheses on the right preserve
the indicated subterm. In particular \(xz(yz)=(xz)(yz)\), not
\(x(z(yz))\), and \(SIx\) contains no \(Ix\)-redex.
The variable-arity reading of combinators is only a convenient notation:
\(S\) still needs three arguments before its binary-tree rule applies.
\end{hbdefinition}

\begin{hbexample}{def:p12-klop1987-combinators}{Cycles, compiled combinators, and equality tests}
The source's cycle is
\(SII(SII)\to I(SII)(I(SII))\to SII(I(SII))\to SII(SII)\).
Figure 7 depicts its additional reductions; the displayed cycle alone
already refutes SN. With \(B=S(KS)K\), the source gives
\[
 Bxyz=S(KS)Kxyz\to KSx(Kx)yz\to S(Kx)yz
 \to Kxz(yz)\to x(yz).
\]
It assigns \(C=S(BBS)(KK)\), with \(Cxyz\twoheadrightarrow xzy\).
The other SKIM definitions are referred to Barendregt and Hindley--Seldin.
Table 5 adds \(Cxyz\to xzy\), \(Bxyz\to x(yz)\),
\(Yx\to x(Yx)\), \(Uz(Pxy)\to zxy\),
\(\mathsf{cond}\ \mathsf{true}\ xy\to x\),
\(\mathsf{cond}\ \mathsf{false}\ xy\to y\), and rules for addition,
multiplication, and equality of numeral constants. The equality instances
\(\mathsf{eq}\ \underline n\ \underline n\to\mathsf{true}\) and
\(\mathsf{eq}\ \underline n\ \underline m\to\mathsf{false}\)
for \(n\ne m\) are ground rule schemata, not a repeated-variable
pattern or a conditional-TRS extension.

Table 6 adjoins a genuine binary operator with \(D(x,x)\to E\).
Table 7 instead adjoins a constant with \(Dxx\to E\), namely
\(\operatorname{Ap}(\operatorname{Ap}(D,x),x)\to E\).
They are not the same signature extension: the latter shares the existing
application operator, whereas the former does not.
\end{hbexample}

\begin{hbdefinition}{def:p12-klop1987-weak-categorical}{The source's Weak Categorical Combinatory Logic}
Table 8, p.15, prints its computational equations with equality signs.
Their left-to-right computational orientation is
\begin{align*}
 \mathsf{Id}\,x&\to x,&(x\circ y)z&\to x(yz),\\
 \mathsf{Fst}(x,y)&\to x,&\mathsf{Snd}(x,y)&\to y,\\
 \langle x,y\rangle z&\to(xz,yz),&\mathsf{App}(x,y)&\to xy,\\
 \Lambda(x)yz&\to x(y,z).
\end{align*}
Here \(\mathsf{Id},\mathsf{Fst},\mathsf{Snd},\mathsf{App}\)
are constants; \(\circ\), the constructor \(\langle\ ,\ \rangle\),
and the constructor \((\ ,\ )\) are binary operators, and
\(\Lambda\) is unary. All the displayed application is still binary
\(\operatorname{Ap}\) in full notation. In particular
\(\mathsf{App}\) is not that underlying application operator, and
\(\mathsf{Fst},\mathsf{Snd}\) are not binary operators despite
their use on an ordered-pair argument. The source attributes its
implementation context to Curien (1986).
\end{hbdefinition}

\begin{hbtheorem}{thm:p12-klop1987-one-rule-puzzle}{proved=annex:p12-klop1987-one-rule-puzzle}{Akkerman's one-rule WN but non-SN answer}
With distinct constants \(a,b\) and binary \(F\), let
\[
 R=\{F(a,F(x,y))\to F(x,F(x,F(b,b)))\}.
\]
Every finite term over this signature has a normal form, but
\(F(a,F(a,F(b,b)))\) has an infinite reduction.
This is the answer to Puzzle 2.1.6, attributed to G. J. Akkerman on p.35.
The normal-form existence claim here has a separate full argument; the
single displayed successful reduction does not by itself establish WN.
\end{hbtheorem}

\begin{hbdefinition}{def:p12-klop1987-direct-sum}{Direct sum, rather than arbitrary union}
For TRSs with disjoint function and constant signatures, \(R_1\oplus R_2\)
is their ordinary union; otherwise rename their nonvariable alphabets
apart before taking the union. Variables can be common.
The report attributes to Toyama (1987) the equivalence
\(\mathrm{CR}(R_1\oplus R_2)\Longleftrightarrow
\mathrm{CR}(R_1)\mathbin{\&}\mathrm{CR}(R_2)\).
This is a modularity theorem, not a claim that arbitrary sharing of symbols
preserves confluence.
The full mixed-term argument is supplied in
Theorem~\ref{thm:p12-klop1987-ds-confluence}, with nonlinear patterns,
duplication and collapse retained and without a termination hypothesis.
\end{hbdefinition}

\begin{hbexample}{def:p12-klop1987-modularity-failures}{The report's two failures of termination modularity}
Let \(R_1=\{f(0,1,x)\to f(x,x,x)\}\) and
\(R_2=\{\mathsf{or}(x,y)\to x,\mathsf{or}(x,y)\to y\}\).
Both terminate separately, but, with \(q=\mathsf{or}(0,1)\),
\(f(q,q,q)\to f(0,q,q)\to f(0,1,q)\to f(q,q,q)\).
The separate factor-termination proofs and this exact cycle are given in
Theorem~\ref{thm:p12-klop1987-ds-first-counter}.

The stronger example has two \(F\)-rules
\(F(4,5,6,x)\to F(x,x,x,x)\), \(F(x,y,z,w)\to7\), and
the nine constant rules
\(1\to4,1\to5,2\to4,2\to6,3\to5,3\to6,4\to7,5\to7,6\to7\).
Its disjoint second system is
\(G(x,x,y)\to x,G(x,y,x)\to x,G(y,x,x)\to x\).
Both are complete, as proved for these exact rules in
Theorem~\ref{thm:p12-klop1987-ds-complete-counter}. For
\(q=G(1,2,3)\), their sum cycles:
\begin{align*}
 F(q,q,q,q)&\twoheadrightarrow
 F(G(4,4,3),G(5,2,5),G(1,6,6),q)\\
 &\twoheadrightarrow F(4,5,6,q)\to F(q,q,q,q).
\end{align*}
The source additionally corrects Ganzinger--Giegerich's p.23 diagram:
the arrow \(b\to k(0)\) must be deleted there, since otherwise
\(f(b,b,b,g(0,b))\) reduces to itself. This is a source-reported
correction to another presentation, not an extra rule in the present example.
\end{hbexample}

The disjoint binary equality-test extension of CL is confluent by
Theorem~\ref{thm:p12-klop1987-ds-applications}, applying the full
direct-sum theorem rather than the regular-confluence theorem.
The applicative equality-test union is reported nonconfluent,
with its argument referred to Klop (1980). The further conditional-selector
\emph{signature} example has the three ordinary rules
\(\mathsf{if}(\mathsf{true},x,y)\to x\),
\(\mathsf{if}(\mathsf{false},x,y)\to y\),
\(\mathsf{if}(z,x,x)\to x\). The report's use of the word ``if''
here denotes a ternary function symbol, not side conditions on rewriting.

% Source-mapped expansion adopted from the released direct-sum fragment.
\subsection{The direct-sum theorem on genuinely mixed terms}
\label{sec:p12-klop1987-direct-sum-proof}

The 1987 tutorial's Theorem~2.3.2 concerns ordinary first-order TRSs
whose \emph{whole nonvariable signatures} are disjoint, not merely
different rule names or different head constants
\parencite[pp.15--17]{Klop1987CWITutorial}. The following full
mixed-term proof reconstructs the preserved-layer and deletion-depth
mechanism of Toyama's separately identified author thesis
\parencite[\S2.2, pp.19--44]{Toyama1990DirectSumProof}.
This later proof source is not treated as a byte-identical copy of
the cited 1987 publisher article or as the 1992 handbook chapter.

\begin{hbdefinition}{def:p12-klop1987-ds-signatures}{The exact unrestricted first-order regime}
Let \(\Sigma_1\cap\Sigma_2=\varnothing\), including constants.
Each operator has a fixed finite arity. Variables belong to one
common infinite supply \(V\). Each rule of \(R_i\) has a
nonvariable left side over \(\Sigma_i\), and its right-side
variables occur on its left. Terms are finite free terms, and
rewriting is closed under arbitrary first-order substitutions and
contexts. No left-linearity, nonduplication, noncollapsing condition
or termination assumption is imposed. The direct sum is
\(R=R_1\oplus R_2\) on \(T(\Sigma_1\cup\Sigma_2,V)\).
If signatures initially share a name, take renamed isomorphic
copies before forming this sum; an ordinary shared-signature union
is a different operation.
For the forward proof assume CR of both components on \emph{all}
their terms, not just ground terms.
\end{hbdefinition}

\begin{hbdefinition}{def:p12-klop1987-ds-layers}{Caps, principal aliens, rank and preservation}
The colour of a nonvariable root is its signature index. At a root
of colour \(a\), the maximal connected \(\Sigma_a\)-cap is
obtained by cutting every first opposite-colour subterm on each
path. These cut subterms are its principal aliens. Variables are
neutral leaves. Write \(t=C[A_1,\ldots,A_k]\) for this occurrence
decomposition; repeated values still have separate occurrence holes.
Its layer rank is 1 for a pure one-signature term or variable, and
\(1+\max_i\operatorname{rank}(A_i)\) for a mixed term.
The special subterms are the root term and the recursively obtained
principal aliens.

Pure one-signature terms and variables are called preserved.
A mixed term of root colour \(a\) is preserved when every reduct
retains a root in \(\Sigma_a\), and all its principal aliens are
preserved. The universal root condition is a mathematical predicate,
not an effective normalizing oracle. It is used below only for an
existence proof; terms failing it are handled by a separate finite
deletion construction.

For the cap proof also use a frontier presentation \(t=H[\mathbf A]\)
that replaces variable leaves by holes as well as principal aliens.
Thus \(H\) contains only operators of one fixed signature and
holes. An empty cap may be a single hole. This presentation is
explicitly different from the principal-alien tuple used for layer rank.
\end{hbdefinition}

\begin{hblemma}{thm:p12-klop1987-ds-invariants}{proved=annex:p12-klop1987-ds-invariants}{Layer invariants without termination}
Rewriting cannot increase layer rank. Preserved terms remain
preserved under reduction. A preserved alien whose initial root
is outside \(\Sigma_a\) never acquires a \(\Sigma_a\)-root;
a pure alien may become a neutral variable, which also remains
outside that signature. These assertions do not bound reduction
length or term size.
\end{hblemma}

\begin{hbdefinition}{def:p12-klop1987-ds-coherence}{The equality condition required by nonlinear patterns}
For equal-length finite tuples write \(\mathbf A\propto\mathbf B\)
if \(A_i=A_j\Rightarrow B_i=B_j\) for all indices. Equality
is literal finite-term equality. Tuple reduction is componentwise.
Coherence does not require distinct inputs to remain distinct.
For a term \(u\), CR at \(u\) means that any two finite
reductions from it join; it is inherited by its descendants.
\end{hbdefinition}

\begin{hblemma}{thm:p12-klop1987-ds-coherence}{proved=annex:p12-klop1987-ds-coherence}{Finite common frontiers and exact cap-step transport}
If each member of a finite frontier tuple has CR, finitely many
chosen descendant tuples have a further common descendant tuple
coherent with each of them and with the original equalities.
A monochromatic cap rewrite whose pattern cannot enter the frontier
values transports to a coherently updated tuple, with the same
rule and the same copied/erased index occurrences on its right.
This covers nonlinear left sides, duplicating right sides and
collapsing right sides.
\end{hblemma}

\begin{hbexample}{def:p12-klop1987-ds-not-commutation}{Why the original component relations need not commute}
Take \(R_1=\{D(x,x)\to E\}\) and \(R_2=\{a\to b\}\),
with disjoint signatures. From \(D(a,a)\), the first component
gives \(E\), while the second gives \(D(b,a)\). No pure
\(R_1\)-path from \(D(b,a)\) joins a pure \(R_2\)-path
from \(E\): neither endpoint admits the needed component step.
The sum joins by \(D(b,a)\to_2D(b,b)\to_1E\).
The tuple update \((a,a)\to(b,a)\) was not coherent; its
further update to \((b,b)\) is. The proof therefore repairs
frontiers, rather than asserting commutation of arbitrary unions.
\end{hbexample}

\begin{hblemma}{thm:p12-klop1987-ds-preserved}{proved=annex:p12-klop1987-ds-preserved}{Preserved caps and their collapsing boundary}
A term consisting of a fixed-colour cap whose frontier aliens are
preserved and have CR itself has CR. Its outer cap may collapse
to one alien. Consequently every preserved term has CR, by
induction on layer rank. The cap statement also proves CR of a
nonpreserved term all of whose principal aliens are preserved.
\end{hblemma}

\begin{hbdefinition}{def:p12-klop1987-ds-depth}{Minimal nonpreserved layers and parallel deletion}
A deletion occurrence is a nonpreserved special subterm all of
whose principal aliens are preserved. Choose a finite reduction
of it to the first term whose root is outside its initial signature.
This root-changing reduct is preserved. A deletion replaces the
occurrence by that reduct, using an actual finite \(R\)-path.
Deletion occurrences are disjoint, because none lies below a
preserved principal alien of another.

Let \(d(t)=0\) for a preserved term; \(d(t)=1\) when \(t\)
itself is a deletion occurrence; otherwise, for a nonpreserved
mixed \(t=C[A_1,\ldots,A_k]\), set
\(d(t)=1+\max_i d(A_i)\). This is the maximum layer depth
of its minimal nonpreserved occurrences. A parallel deletion
\(t\Rightarrow_{\rm d}u\) deletes all occurrences at that
maximum depth, with independently chosen finite witnesses.
No computable selector for this semantic classification is claimed.
\end{hbdefinition}

\begin{hblemma}{thm:p12-klop1987-ds-depth}{proved=annex:p12-klop1987-ds-depth}{The actual decreasing depth}
The depth is a finite natural number. Ordinary rewriting cannot
increase it. Every nonpreserved term has a parallel deletion;
such a deletion is a finite ordinary reduction and strictly
decreases \(d\). Every deletion occurrence has CR by the cap
lemma. This does not assert termination of ordinary rewriting.
\end{hblemma}

\begin{hblemma}{thm:p12-klop1987-ds-diagrams}{proved=annex:p12-klop1987-ds-diagrams}{All ordinary-step/parallel-deletion positions}
Given \(t\to u\) and \(t\Rightarrow_{\rm d}v\), with
\(d(t)>0\), the following joins exist.
If \(d(u)=d(t)\), one can choose
\(u\Rightarrow_{\rm d}w\), with \(v\to^*z,w\to^*z\)
and \(d(z)<d(t)\). If \(d(u)<d(t)\), one can choose
\(v\to^*z,u\to^*z\) with that same strict bound.
The proof treats disjoint steps, a step inside a selected layer,
and a step above selected layers. In the last case coherent
joining is essential when copies reduce differently.
\end{hblemma}

\begin{hbtheorem}{thm:p12-klop1987-ds-confluence}{proved=annex:p12-klop1987-ds-confluence}{The tutorial's unrestricted direct-sum theorem}
For the exact regime above,
\[
 \mathrm{CR}(R_1\oplus R_2)
 \quad\Longleftrightarrow\quad
 \mathrm{CR}(R_1)\ \mathbin{\&}\ \mathrm{CR}(R_2).
\]
This includes mixed terms, nonlinear patterns, duplication and
collapse. It does not assume either component terminates, and
does not replace the sum by a shared-signature union.
\end{hbtheorem}

\begin{hbtheorem}{thm:p12-klop1987-ds-first-counter}{proved=annex:p12-klop1987-ds-first-counter}{The source's first termination failure}
Let \(R_1=\{f(0,1,x)\to f(x,x,x)\}\) and
\(R_2=\{\mathsf{or}(x,y)\to x,\mathsf{or}(x,y)\to y\}\),
over disjoint signatures. Both terminate. With
\(q=\mathsf{or}(0,1)\), their sum has the cycle
\[
 f(q,q,q)\to f(0,q,q)\to f(0,1,q)\to f(q,q,q).
\]
Here \(R_2\) is not CR. This is the particular source counterexample,
not an argument that two complete factors can never be combined.
\end{hbtheorem}

\begin{hbdefinition}{def:p12-klop1987-ds-complete-counter}{The exact eleven-rule/four-argument source factor}
In the next source example the glyphs 1 through 7 are distinct
nullary symbols, not numbers in a semantic ordering. The first
factor has the two rules
\(F(4,5,6,x)\to F(x,x,x,x)\),
\(F(x,y,z,w)\to7\), and all nine constant rules
\[
1\to4,\ 1\to5,\ 2\to4,\ 2\to6,\ 3\to5,\ 3\to6,
\quad4\to7,\ 5\to7,\ 6\to7.
\]
The disjoint second factor has exactly
\(G(x,x,y)\to x\), \(G(x,y,x)\to x\),
\(G(y,x,x)\to x\).
These are the whole source graph and its three majority rules.
\end{hbdefinition}

\begin{hbtheorem}{thm:p12-klop1987-ds-complete-counter}{proved=annex:p12-klop1987-ds-complete-counter}{Completeness is also not modular without extra hypotheses}
Both factors in the preceding definition are SN and CR. Their
sum is CR by the full modularity theorem, but is not SN.
Putting \(q=G(1,2,3)\), the source cycle is
\begin{align*}
 F(q,q,q,q)&\to^*
 F(G(4,4,3),G(5,2,5),G(1,6,6),q)\\
 &\to^*F(4,5,6,q)\to F(q,q,q,q).
\end{align*}
The source attributes this construction to Barendregt and Klop,
adapted by Toyama. It refutes modularity of SN and of CR plus SN;
it does not refute confluence modularity.
\end{hbtheorem}

\begin{hbtheorem}{thm:p12-klop1987-ds-applications}{proved=annex:p12-klop1987-ds-applications}{The binary equality test and selector applications}
The one-rule binary factor \(D(x,x)\to E\) is SN and CR.
The factor with \(\mathsf{if}(\mathsf{true},x,y)\to x\),
\(\mathsf{if}(\mathsf{false},x,y)\to y\),
\(\mathsf{if}(z,x,x)\to x\) is also SN and CR.
Their separate direct sums with CL are therefore CR. The nonlinear
test and selector do not become regular by being disjoint factors.
Here \(\mathsf{if}\) is an ordinary ternary symbol, not a condition
attached to a rule.
\end{hbtheorem}

The report separately contrasts the applicative test \(Dxx\to E\).
Its restored left side is
\(\operatorname{Ap}(\operatorname{Ap}(D,x),x)\), so that
extension shares the application operator with CL. It is not this
direct sum; the source reports it nonconfluent and refers its proof
to Klop(1980). That separate nonconfluence witness is not supplied
by the binary-factor proof or by a counterexample to an unrelated union.
The report also mentions reduced/irreducible variants in Toyama's
1986 preprint and Ganzinger--Giegerich(1987), and reports the latter's
extra \(b\to k(0)\) arrow as producing a self-reduction of
\(f(b,b,b,g(0,b))\). These cited variants retain their own rule
sets and proof standing, rather than being added to the eleven-rule
factor above.


\subsection{Semantics and the source's finite decision questions}

\begin{hbtheorem}{thm:p12-klop1987-semantic-un}{proved=annex:p12-klop1987-semantic-un}{The source's semantic test for UN}
Suppose an algebra of the same signature satisfies every rule equation
under every variable assignment, and distinct syntactic normal forms are
not equal in that algebra. Then the TRS has UN. This is Theorem 2.4.1,
p.17, and applies to open terms as well as ground terms.
\end{hbtheorem}

Section 2.5 fixes finite signatures and finite rule sets for its decision
questions. It reports undecidability of confluence and termination in
general, but decidability of each for ground TRSs; the all-terms/ground-terms
restriction is kept separate from ground rules. It attributes two-rule
termination undecidability to Dershowitz and calls the one-rule question
open \emph{at the date of the report}. It also reports undecidability,
within CL, of conversion and existence of finite normal forms or infinite
reductions. Its last report is equality of the ground-normal-form sets of
two left-linear TRSs, attributed to Wiedijk (1987), with removal of
left-linearity described only as work in progress. None of these dated
sentences asserts the corresponding state of research in 2026.

\subsection{The marked-tree termination route}

\begin{hbdefinition}{def:p12-klop1987-star-order}{Commutative trees and the five marker operations}
The report's \(\mathcal T\) consists of finite unordered rooted trees
labelled in \(\mathbb N\); \(\mathcal T^*\) permits one marker on a
node. Write \(n(\vec t)\) for the children and \(t^*\) for a root
marker. Its reduction permits:
placing a root marker \(n(\vec t)\to n^*(\vec t)\);
copying under a smaller label
\(n^*(\vec t)\to m(n^*(\vec t),\ldots,n^*(\vec t))\) when \(m<n\),
with any finite number of copies, including zero;
pushing a marker into a chosen child,
\(n^*(s,\vec t)\to n(s^*,\ldots,s^*,\vec t)\), again with zero
or more copies, with the chosen child's root initially unmarked; selecting a child
\(n^*(t_1,\ldots,t_k)\to t_i\) for \(k>0\); and performing
any such operation within a tree context.
Commutativity concerns sibling order, not associativity of operators.
\end{hbdefinition}

The full marked relation is not terminating: copying may leave a copy of
the marked input available for repetition. The source's Theorem 3.5 instead
concerns positive reductions \emph{between unmarked trees}. Its dependency
is Kruskal's theorem for natural-number-labelled finite trees: every
infinite sequence has an earlier tree embedded in a later one by the
report's ancestry-order convention, with labels nondecreasing. This is
the convention specified at \cref{def:p12-klop1987-mt-embedding}, not an
unstated least-common-ancestor preservation requirement. The argument combines
this with acyclicity of the marker order and its compatibility with proper
embedding. Embedding includes the identical tree: its reduction comparison
is reflexive in that case and positive for \emph{proper} embedding.
The positive relation on unmarked trees has no self-cycle.

In Table 9 the five formula-rewriting rules are double-negation removal,
the two De Morgan rules, and the two distributive rules sending conjunction
over disjunction. The assignment \(\lor\mapsto1,\land\mapsto2,
\neg\mapsto3\) sends them to decreasing marker reductions. For instance,
\[
3(1(t,s))\to3^*(1(t,s))
\to2(3^*(1(t,s)),3^*(1(t,s)))
\to^*2(3(t),3(s)).
\]
The last segment pushes each marker onto the \(1\)-child and selects
the appropriate argument. The other four comparisons are
\(3(3(t))\to^+t\), \(3(2(t,s))\to^+1(3(t),3(s))\),
\(2(t,1(s,r))\to^+1(2(t,s),2(t,r))\), and its symmetric variant.
The mathematical content is termination of these syntactic transformations,
not completeness of a logical calculus.

Remark 3.11 gives three boundaries. An embedded left side in a right side,
as in \(f(s(x))\to g(s(x),f(p(s(x))))\), obstructs this method.
Taken as an isolated one-rule system, with no rules for \(p\) or \(s\),
that example has the different active-origin termination proof of
Theorem~\ref{thm:p12-klop1987-pa-puel-example}. The companion-rule
failure in Lemma~\ref{thm:p12-klop1987-pa-puel-extension-failure}
shows why this does not establish arbitrary extensions or Puel's
general unavoidable-set theorem.
The process-algebra rules can reintroduce a heavier operator. The source
proposes norm-indexed copies using symbol lengths, but automatic fresh
reannotation after PA duplication need not respect an enclosing context:
parallel expansion can raise an outer left-merge index from 4 to 8, as
Lemma~\ref{thm:p12-klop1987-pa-reannotation-failure} proves.
The actual three source PA rules have complete contextual certificates
using the distinct max/sum annotation and retained initial-tag budgets
of Theorems~\ref{thm:p12-klop1987-pa-termination} and
\ref{thm:p12-klop1987-pa-budget-lift}. Their exact ordinary
three-rule and seven-rule scopes do not include the larger ACP system
or implicit rewriting modulo choice equations.
The simplified example is \(g(x,y)\to h(x,y)\),
\(h(f(x),y)\to f(g(x,y))\), with
\(g_n>h_n>f\) and \(h_{n+1}>g_n\).
Associativity and the Ackermann rules require a lexicographic, ordered-tree
extension, attributed to Kamin--Levy through Dershowitz (1985).
The Ackermann comparison uses
\[
 A(S(x),S(y))\to A^*(S(x),S(y))
 \to A(S^*(x),A^*(S(x),S(y)))
 \to^* A(x,A(S(x),y)).
\]
This extension preserves fixed arities and left-to-right order; it is not
the commutative-tree order of Definition~\ref{def:p12-klop1987-star-order}.
Its global plain-endpoint well-foundedness is established separately in
Theorem~\ref{thm:p12-klop1987-om-well-founded};
Theorem~\ref{thm:p12-klop1987-om-applications} supplies the positive
paths for all three Ackermann rules and for ordered associativity.

% ED26-SRC-027-501 continuation. Integrator-owned placement in source section3.
% Does not edit or supersede the released core while integration is active.
\subsection{The marked-tree argument at its original hypotheses}
\label{sec:p12-klop1987-marked-tree-proofs}

The following expands the argument in Klop's 1987 report,
\S3, pp.18--22 \parencite{Klop1987CWITutorial}. Its plain objects are
finite, nonempty, rooted trees with natural-number labels and unordered
children. Operator arities are unrestricted in this part of the report.
The stars are temporary control marks, not additional numbers in the
label order. A full tracing proof of the missing acyclicity step is
provided here with the separately attributed argument of
\textcite[\S4.3]{EndrullisKlopOverbeek2021MarkedTrees}; this later proof
dependency is not identified with the 1992 chapter or with a proof
printed in the 1987 tutorial.

\begin{hbdefinition}{def:p12-klop1987-mt-objects}{Unordered tree objects and well-defined marks}
Represent a tree by \(n(t_1,\ldots,t_k)\), with \(n\in\mathbb N\)
and any finite \(k\geq0\), identifying representations by permutations
of siblings at every node. This identification is denoted \(\simeq\)
on representations. Its equivalence classes form \(\mathcal T\).
Allowing each node either its plain label \(n\) or its marked version
\(n^*\) gives \(\mathcal T^*\). Each node has at most one mark.
For a tree with unmarked root \(t=n(\vec u)\), \(t^*=n^*(\vec u)\)
marks that root; it does not recursively mark the children.

The report's operations are putting this root mark; selecting one
child of a marked root; replacing \(n^*(\vec u)\) by
\(m(n^*(\vec u),\ldots,n^*(\vec u))\) when \(m<n\), with any
finite number of copies, including zero; and pushing the mark to a
chosen unmarked-root child \(v\), replacing
\(n^*(v,\vec u)\) by \(n(v^*,\ldots,v^*,\vec u)\), again with
zero or more copies. Thus the last operation removes the parent's mark.
The type of \(v^*\) requires that its root was unmarked; a second mark
on an already marked root is not silently added.
The source reduction is \(\to\), with its stated context closure.
For a proof comparison only, write \(\rightsquigarrow\) for the
same operations closed under \emph{all} finite contexts, including
marked ones. Every source step is a \(\rightsquigarrow\)-step.
This auxiliary enlargement makes no claim that a printed unmarked
context clause and every marked context are the same definition.
\end{hbdefinition}

\begin{hbdefinition}{def:p12-klop1987-mt-embedding}{The report's ancestry-order embedding}
For nodes \(p,q\) of a rooted tree, let \(p\leq_{\rm anc}q\)
mean that \(p\) is an ancestor of \(q\), including equality.
Write \(s\sqsubseteq t\) when there is an injective node map
\(\phi\) satisfying
\[
 p\leq_{\rm anc}q\quad\Longleftrightarrow\quad
 \phi(p)\leq_{\rm anc}\phi(q),\qquad
 \operatorname{label}(p)\leq\operatorname{label}(\phi(p)).
\]
The root may map to a proper descendant of the target root.
The equivalence on the first line makes the map an isomorphism onto
its induced ancestry-ordered image, as specified in the report.
It need not preserve the least common ancestor of two incomparable
nodes. This exact distinction is material to its Figure~10.

For the well-quasi-order proof, temporarily give children a planar
order and use a stronger relation \(\preccurlyeq_{\rm pl}\).
It permits diving into a target child; or coupling roots with
nondecreasing labels and embedding the source child list into a
subsequence of the target child list, using the same relation on
each paired child. Coupling places different source children in
different immediate target branches. Every such planar embedding,
after forgetting sibling order, is an ancestry-order embedding above.
\end{hbdefinition}

\begin{hblemma}{thm:p12-klop1987-mt-embedding}{proved=annex:p12-klop1987-mt-embedding}{Well-defined embedding conventions and their distinct strength}
The source relation \(\sqsubseteq\) is a reflexive, transitive,
antisymmetric relation on unordered labelled-tree classes. The planar
relation is a preorder, and every planar embedding induces a source
embedding. Requiring least-common-ancestor preservation is stronger
than the ancestry-order convention, not an alternative notation for it.
\end{hblemma}

\begin{hbexample}{def:p12-klop1987-mt-figures8-9}{The actual marked trees and the copying example}
Figure~8's plain tree is \(3(5,7(9),8(0(1,5)))\), and its marked
version is \(3^*(5,7(9^*),8^*(0(1,5)))\). The marks are separate
node bits. For Figure~9 put \(a=5(6,7(8))\). Its construction is
\begin{align*}
 a&\to a^*\to4(a^*,a^*,a^*)\\
 &\to^*4(6,5(6,7^*(8)),4(a^*,a^*))\\
 &\to^*4(6,5(6,8),4(6,7(8))).
\end{align*}
The middle segment selects 6 in the first copy, pushes once into 7(8)
in the second, and copies the third twice under 4. The last segment
selects 8 from 7*, then 6 and 7(8) from the last marked 5-copies.
The displayed finite construction ends unmarked. Repeatedly copying
a still marked 5-subtree instead gives an infinite auxiliary reduction.
\end{hbexample}

\begin{hbexample}{def:p12-klop1987-mt-figure10}{The actual 1987 embedding example and the later variant}
The report has
\[
 2(9,7(0,4))\sqsubseteq1(3(8(0(5,1)),9,5(9)),2).
\]
Map its root 2 to the 3-node; its 9-child to that node's 9-child;
its 7-child to the 8-node; and the 0- and 4-children to the 1- and 5-leaves
below that 8-node, respectively. Each label inequality holds, and
ancestry is preserved and reflected. The last two images have the 0-node
as their least common ancestor, not the image 8 of their source parent 7.
Thus this example must not be used to assert meet preservation.
The later primary paper uses instead an 8 in place of that internal 0
and an explicitly meet-preserving embedding
\parencite[Example4.3]{EndrullisKlopOverbeek2021MarkedTrees}.
The stronger planar relation below is used to prove a well-quasi-order
for the report's weaker relation, without collapsing these variants.
\end{hbexample}

\begin{hbdefinition}{def:p12-klop1987-mt-wqo}{The required quasi-order property}
A quasi-order is reflexive and transitive. It is a well-quasi-order
(wqo) if every infinite sequence has indices \(i<j\) with its
\(i\)-th member below its \(j\)-th member. A bad sequence lacks
such a pair. Finite words are ordered by increasing subsequence
embedding, allowing deletion of letters and replacement by smaller
letters. Empty words are included.
All alphabets used in the following constructions have a countable
coding: natural labels, finite trees over those labels, and their
finite lists. Minimal bad sequences are chosen by length or number
of nodes, with a fixed enumeration breaking ties; this is a classical
existence argument, not an algorithm deciding whether a bad extension
exists.
\end{hbdefinition}

\begin{hblemma}{thm:p12-klop1987-mt-words}{proved=annex:p12-klop1987-mt-words}{The word and product dependencies}
Every infinite sequence in a wqo has an infinite nondecreasing
subsequence. Finite words over a countably coded wqo form a wqo under
subsequence embedding, and the product of two wqos is a wqo under
componentwise comparison. The natural numbers with their usual order
satisfy these hypotheses.
\end{hblemma}

\begin{hbtheorem}{thm:p12-klop1987-mt-kruskal}{proved=annex:p12-klop1987-mt-kruskal}{The full labelled-tree wqo needed in the report}
Finite planar natural-number-labelled trees are wqo under
\(\preccurlyeq_{\rm pl}\). Consequently the report's unordered
trees are wqo under \(\sqsubseteq\). Labels, depths and branching
degrees need not have one common finite bound. This proves the
natural-number version of the dependency invoked at Theorem~3.8,
not merely a fixed finite-ranked-alphabet special case.
\end{hbtheorem}

\begin{hblemma}{thm:p12-klop1987-mt-embedding-reduction}{proved=annex:p12-klop1987-mt-embedding-reduction}{Embedding is simulated in the report's marker relation}
If \(s\sqsubseteq t\), then \(t\to^*s\). If the tree classes
are different, the reduction is positive. The identical case uses
zero steps; the unrestricted assertion \(t\to^+t\) is not inferred
from reflexive embedding. The simulation handles distinct source
branches whose image roots share an initial target stem.
\end{hblemma}

\begin{hblemma}{thm:p12-klop1987-mt-tracing}{proved=annex:p12-klop1987-mt-tracing}{The occurrence-level marker mechanism}
In the all-context auxiliary game, each output node has one input
ancestor. Its numerical label does not exceed that ancestor's label;
a new lower root introduced by copying is a strict decrease on that
lineage. A tracked subtree has a projected reduction using at most
one step per original step. Finite sibling permutations can be
postponed, so a cycle on tree classes would give a marker reduction
on planar representations ending in a permutation of its input.
\end{hblemma}

\begin{hbtheorem}{thm:p12-klop1987-mt-acyclic}{proved=annex:p12-klop1987-mt-acyclic}{Acyclicity on plain trees, including the zero-copy cases}
There is no positive \(\rightsquigarrow\)-reduction between an
unmarked tree and itself. Hence the report's source relation has
no such cycle. The assertion is not acyclicity or termination of
the whole marked auxiliary system.
\end{hbtheorem}

\begin{hbtheorem}{thm:p12-klop1987-mt-well-founded}{proved=annex:p12-klop1987-mt-well-founded}{The report's unmarked strict order is well-founded}
The relation \(\to^+\) restricted to pairs in \(\mathcal T\)
is a strict, transitive, well-founded order. In particular there is
no infinite sequence of unmarked trees joined successively by
positive finite marker reductions, as asserted in Theorem~3.5.
The proof uses the established wqo, embedding simulation and
acyclicity; none is replaced by its name.
\end{hbtheorem}

The marked universe is genuinely different. For example,
\(1^*()\to0(1^*())\to0(0(1^*()))\to\cdots\)
uses repeated copying with one copy and never becomes entirely
unmarked. It witnesses failure of SN on \(\mathcal T^*\), exactly
the distinction made before the report's Theorem~3.5.

\begin{hbtheorem}{thm:p12-klop1987-mt-boolean}{proved=annex:p12-klop1987-mt-boolean}{All five source Boolean-formula rewrite comparisons}
The five rules of source Table~9 terminate: double-negation removal,
both De Morgan transformations, and both distributive transformations
of conjunction over disjunction. Their positive marker simulations
use \(\lor\mapsto1\), \(\land\mapsto2\), \(\neg\mapsto3\).
This is termination of these syntactic rules, not confluence or
completeness of a logical proof calculus.
\end{hbtheorem}

\begin{hbtheorem}{thm:p12-klop1987-mt-indexed}{proved=annex:p12-klop1987-mt-indexed}{The norm-indexed source pair, with the context dependency checked}
For the two rules \(g(x,y)\to h(x,y)\) and
\(h(f(x),y)\to f(g(x,y))\), the source's norm-indexed construction
proves termination. Annotate a binary \(g(s,r)\) or \(h(s,r)\)
by \(n=|s|+|r|\), with bars counting nodes, and take labels
\(g_n=2n\), \(h_n=2n-1\), \(f=0\).
Then \(g_n>h_n>f\) and \(h_{n+1}>g_n\).
Both original rules preserve node count, so ancestor annotations
stay unchanged when a rule is used in an arbitrary context.
This establishes the concrete two-rule construction; it does not
silently extend node-count preservation to the full process-algebra
rules that also duplicate arguments under a sum.
\end{hbtheorem}

\begin{hbcorollary}{thm:p12-klop1987-mt-obstruction}{proved=annex:p12-klop1987-mt-obstruction}{The source's embedded-left-side obstruction}
An ordinary rule whose left side embeds in its right side cannot have
a positive descending simulation in the same acyclic marker order.
The source example \(f(s(x))\to g(s(x),f(p(s(x))))\) has that
obstruction whenever each operator's label is unchanged between its
two occurrences.
\end{hbcorollary}

\begin{hbdefinition}{def:p12-klop1987-mt-ordered-extension}{The separately enlarged ordered marker relation}
Remark~3.11(iii) changes the objects to ordered trees with fixed
operator arities. Copying to a lower label must give that lower
operator its prescribed number of arguments. Pushing changes one
chosen argument \(t_i\) to \(t_i^*\), without deleting or copying
that argument. For \(n\) of arity \(k\geq1\), the added
left-argument operation is
\[
 n^*(t_1,\ldots,t_k)\longrightarrow
 n(t_1^*,n^*(t_1,\ldots,t_k),\ldots,n^*(t_1,\ldots,t_k)),
\]
with exactly \(k-1\) full copies after the first argument.
The marker on \(t_1\) is defined at an unmarked root as above.
The report gives the Ackermann comparison
\begin{align*}
 A(S(x),S(y))&\longrightarrow A^*(S(x),S(y))\\
 &\longrightarrow A(S^*(x),A^*(S(x),S(y)))\\
 &\longrightarrow A(x,A^*(S(x),S(y)))\\
 &\longrightarrow A(x,A(S(x),S^*(y)))\\
 &\longrightarrow A(x,A(S(x),y)).
\end{align*}
Its source rules are \(A(0,x)\to S(x)\),
\(A(S(x),0)\to A(x,S(0))\), and
\(A(S(x),S(y))\to A(x,A(S(x),y))\).
This displayed derivation belongs to the enlarged ordered relation,
not to the unranked commutative relation proved well-founded above.
\end{hbdefinition}

\begin{proofboundary}{def:p12-klop1987-mt-boundaries}{Separate mechanisms and separate source standing}
The main proof concerns the unranked natural-labelled source tree
order. Adding the left-argument operation changes the reduction
relation, so inclusion in the proved four-operation game is unavailable.
The global well-foundedness of that enlarged relation is instead
established by the separate lower-set interpretation in
Theorem~\ref{thm:p12-klop1987-om-well-founded}; its strict conclusion
concerns plain endpoints, not the whole marked universe. An Ackermann
marker comparison alone would not supply that argument. Likewise, the displayed
rule whose left side embeds in its right side requires a different
method: \(f(s(x))\to g(s(x),f(p(s(x))))\) cannot strictly descend
in an acyclic order also containing reverse embedding.
The report refers that latter extension to Puel(1986), and refers
the ordered lexicographic extension to Kamin--Levy through Dershowitz.
The ordered-extension proof does not discharge the distinct general
Puel extension. The actual three ordinary PA rules instead have the
separate contextual max/sum and retained-budget proofs in
Theorems~\ref{thm:p12-klop1987-pa-termination} and
\ref{thm:p12-klop1987-pa-budget-lift}, at their declared three-rule
and seven-rule scopes. The larger ACP communication/silent-action
variants and Puel's general unavoidable-set construction remain outside
these proofs; neither is established by the isolated example's
active-origin argument.
\end{proofboundary}

% Source-mapped PA/Puel expansion; seven original reciprocal proof pairs.
\subsection{Norm-indexed PA: the contextual obligation and its repair}
\label{sec:p12-klop1987-pa-norm-proof}

Remark~3.11(ii), printed p.21, gives the three ordinary rules
\begin{align*}
 x\parallel y&\to(x\mathbin{\lhd}y)+(y\mathbin{\lhd}x),\\
 (x+y)\mathbin{\lhd}z&\to(x\mathbin{\lhd}z)+(y\mathbin{\lhd}z),\\
 (a\cdot x)\mathbin{\lhd}y&\to a\cdot(x\parallel y).
\end{align*}
Here \(\lhd\) denotes the source's binary left-merge glyph, not
an inequality; \(a\) ranges over atomic-action constants. The
unindexed precedence \(\parallel>\lhd>\cdot,+\) does not handle
the last rule, which recreates the heavier operator. The report
introduces indexed operators and illustrates the method on its
node-count-preserving \(g/h/f\) pair. That pair's proof does not
by itself justify annotations above a duplicating PA step
\parencite[p.~21]{Klop1987CWITutorial}.

The earlier primary CWI report CS-R8403 has the same indexing idea
in Appendix I,13, but defines an additive atomic-length norm, including
addition across \(+\), rather than counting every syntax node
\parencite[pp.~30--34]{BergstraKlop1984PAIndexingReport}.
These are different norms. The construction below both exhibits the
contextual reannotation failure and supplies a complete repaired
max/sum interpretation. It also gives a budget interpretation that
can retain the tutorial's original \emph{initial} syntax-length tags.
Neither repair is silently identified with a full proof printed in
either source.

\begin{hbdefinition}{def:p12-klop1987-pa-regime}{The exact ordinary PA rewrite systems being proved}
Let \(A\) be a nonempty set of distinct nullary atomic-action symbols
and \(V\) an infinite variable supply disjoint from it. Finite free
terms have those leaves and four binary symbols
\(+,\cdot,\parallel,\lhd\). Let \(R_3\) consist exactly of
the three displayed rules, including the action-indexed third schema.
Let \(R_7\) additionally contain
\[
 a\mathbin{\lhd}x\to a\cdot x,\quad x+x\to x,\quad
 (x+y)\cdot z\to(x\cdot z)+(y\cdot z),\quad
 (x\cdot y)\cdot z\to x\cdot(y\cdot z).
\]
The first additional schema is the atomic left-merge rule; the other
three are the oriented idempotence, distribution and sequential
associativity equations. The primary table calls these CM2, A3,
A4 and A5. Its ACP merge equation also includes a communication
merge summand, which is \emph{not} in the tutorial's PA rule above.
That primary report assumes a finite action alphabet; the proof below
also covers an infinite ambient action set by establishing a finite
alphabet for each individual run, rather than attributing that
extension to the primary statement.

All rewriting here is ordinary rewriting on these free finite terms,
under arbitrary finite-term instances and contexts. Commutativity
and associativity of \(+\) are not directed rules or implicit
matching equations. No communication merge, deadlock, silent action,
encapsulation, abstraction or recursive-definition rule is added.
Thus \(R_7\) is the expressly stated ordinary PA presentation;
it is not all of the primary report's larger ACP system. That
appendix in fact excludes A5 from its commutative-tree argument;
the present A5 case uses ordered left-to-right LPO separately.
\end{hbdefinition}

\begin{hblemma}{thm:p12-klop1987-pa-reannotation-failure}{proved=annex:p12-klop1987-pa-reannotation-failure}{Literal length reannotation does not lift every contextual PA step}
Take \(t=(a\parallel b)\mathbin{\lhd}c\) and perform the first
source rule inside its left argument, giving
\(u=((a\mathbin{\lhd}b)+(b\mathbin{\lhd}a))\mathbin{\lhd}c\).
Annotating every merge in every successive term afresh by the sum
of its arguments' syntax-node counts changes the outer left-merge
index from 4 to 8. With the primary additive atomic-length norm applied
to this two-branch PA rule it changes from 3 to 5. Applying the primary
report's actual three-branch CM1 rule, with its additional communication
merge \(a\mathbin{|}b\), changes that primary norm's outer index
from 3 to 7. In each case the new outer numerical label
exceeds every old numerical label in any order realizing the source's
interleaved indexed precedence. The source marker operations cannot
produce such a larger label. Therefore this automatic reannotation
does not give the required positive marker comparison in an arbitrary
context. This refutes that proposed lift, not termination of PA.
\end{hblemma}

\begin{hbdefinition}{def:p12-klop1987-pa-stable-norm}{A positive max/sum norm and its canonical annotation}
Define a different norm \(\nu\) by
\[
 \nu(a)=\nu(x)=1,\qquad
 \nu(p+q)=\max(\nu(p),\nu(q)),\qquad
 \nu(p\circ q)=\nu(p)+\nu(q)
 \quad(\circ\in\{\cdot,\parallel,\lhd\}).
\]
Thus every action and variable has positive unit weight, and choice
uses maximum, not sum or syntax size. Let \(E(t)\) preserve variables,
actions, \(+\) and \(\cdot\), recursively annotate the children,
and replace \(p\parallel q\) and \(p\lhd q\) by
\(P_{\nu(p)+\nu(q)}(E(p),E(q))\) and
\(L_{\nu(p)+\nu(q)}(E(p),E(q))\), respectively.
All \(P_n,L_n\) have fixed arity two and \(n\geq2\).

Use left-to-right LPO with precedence ordered first by indexed level,
then \(P_n>L_n\) within a level, and all indexed symbols above
\(\cdot>+>a\) for every \(a\in A\). No comparison between
different actions is needed. In particular
\(P_n>L_n\) and \(L_{n+1}>P_n\), retaining the source's
required interleaving. This is an infinite ambient indexed signature;
the run-level finite-signature obligation is proved below.
\end{hbdefinition}

\begin{hblemma}{thm:p12-klop1987-pa-norm}{proved=annex:p12-klop1987-pa-norm}{Invariance for every instance and every ancestor}
The norm is a well-defined positive integer on every finite term;
every subterm has norm at most that of the whole term. Every
\(R_7\)-step preserves it, including under arbitrary substitutions
and contexts. Norm evaluation after a substitution assigns each
variable the positive norm of its substituted term; no unjustified
identity between an indexed annotation and an ordinary substitution
is assumed. A change with equal norm leaves all ancestor merge
indices unchanged, so both endpoints have one common annotated
one-hole context around the changed subtree.
\end{hblemma}

For example, \(E(x\parallel y)=P_2(x,y)\), but substituting
\(a\cdot a\) for \(x\) gives canonical annotation
\(P_3(a\cdot a,y)\), not \(P_2(a\cdot a,y)\).
Consequently the proof checks actual instances first, rather than
asserting that ordinary substitution preserves numerical annotations.

\begin{hblemma}{thm:p12-klop1987-pa-lpo-lift}{proved=annex:p12-klop1987-pa-lpo-lift}{Every actual PA rule instance strictly decreases LPO}
For every ordinary instance \(l\sigma\to r\sigma\) of each
of the seven stated rules,
\(E(l\sigma)>_{\rm lpo}E(r\sigma)\). Consequently every
contextual step \(t\to_{R_7}u\) gives
\(E(t)>_{\rm lpo}E(u)\). The proof supplies all domination
and equal-head lexicographic cases, including both independently
sized branches of distribution and the recreated
\(P_{n-1}\) below a consuming \(L_n\).
\end{hblemma}

\begin{hbtheorem}{thm:p12-klop1987-pa-termination}{proved=annex:p12-klop1987-pa-termination}{The complete ordinary PA termination argument}
Both \(R_3\) and \(R_7\) terminate on all finite open and ground
terms of their declared signatures. In a run starting at \(t\),
every annotated merge index is between 2 and \(N=\nu(t)\).
Only the finitely many action and variable symbols already in \(t\),
the two fixed base operators, and those finitely many indexed operators
are needed throughout that entire run. There is no common global
bound on indices or term size imposed on all inputs.
The finite-ranked LPO well-foundedness proof therefore applies to
the actual run, rather than being invoked over an unverified infinite
signature. The result is SN, not a confluence or PA completeness theorem.
\end{hbtheorem}

\begin{hbdefinition}{def:p12-klop1987-pa-budget}{Retained initial tags as upper budgets}
For the three-rule system \(R_3\), an admissible annotated term
has erasure in the ordinary signature and satisfies
\(n\geq\nu(p)+\nu(q)\) at each \(P_n(p,q)\) or
\(L_n(p,q)\), with \(\nu\) evaluated on erasures.
This is an upper-bound condition, not equality to the current
syntax length. Its lifted rules are
\begin{align*}
 P_n(p,q)&\to L_n(p,q)+L_n(q,p),\\
 L_n(p+q,r)&\to L_n(p,r)+L_n(q,r),\\
 L_n(a\cdot p,q)&\to a\cdot P_{n-1}(p,q).
\end{align*}
They keep all surrounding tags and use the inherited \(n\) in
each distributed branch. The last rule is defined only where
admissibility gives \(n\geq3\).
\end{hbdefinition}

\begin{hbtheorem}{thm:p12-klop1987-pa-budget-lift}{proved=annex:p12-klop1987-pa-budget-lift}{A complete lift retaining the tutorial's initial syntax-length indices}
Every finite ordinary term has an admissible annotation obtained by
the tutorial's initial syntax-length prescription. From any admissible
annotation of an \(R_3\)-source, every actual ordinary step lifts
to the displayed annotated rule in the same unchanged-tag context.
Admissibility persists, and each lifted step strictly decreases the
stated LPO. If an initial index is present, newly created indices
do not exceed the largest initial index; without any merge node,
the term has no \(R_3\)-step. This supplies another complete termination route for the
actual three source rules, without falsely reannotating every later
term by its enlarged syntax length. It does not claim a coherent
upper-budget lift for every nonlinear or larger ACP rule.
\end{hbtheorem}

\subsection{The distinct embedded-left-side example}

The report's Remark~3.11(i) refers the obstruction
\(f(s(x))\to g(s(x),f(p(s(x))))\) to Puel's extension of
Kruskal's theorem. Its left side embeds in its right side by mapping
the first \(f\) to the new inner \(f\), and diving through
the inserted \(p\) before matching \(s\) and its descendants.
Thus a simplification order containing reverse embedding cannot
orient it, as established in
Corollary~\ref{thm:p12-klop1987-mt-obstruction}. That obstruction
is different from PA's re-emerging heavy operator and is not removed
merely by renaming its symbols as merge operators.

\begin{hbtheorem}{thm:p12-klop1987-pa-puel-example}{proved=annex:p12-klop1987-pa-puel-example}{Termination of the actual isolated source example by a different measure}
For the one-rule ordinary TRS
\[
 Q=\{f(s(x))\to g(s(x),f(p(s(x))))\},
\]
with unary \(f,s,p\), binary \(g\), variables and any inert
constants, every finite term terminates. No rule for \(p\) or
\(s\) is present. The proof traces the initially active \(f\)
occurrences: a contracted occurrence disappears, the new \(f\)
above a \(p\)-root is permanently inactive, and only strictly
deeper old active origins are copied. A finite lexicographic residual
count proves termination despite the growth of syntax.
This is a concrete proof of the isolated displayed system, not a
reconstruction of Puel's general unavoidable-set theorem.
\end{hbtheorem}

\begin{hblemma}{thm:p12-klop1987-pa-puel-extension-failure}{proved=annex:p12-klop1987-pa-puel-extension-failure}{Why the isolated proof does not settle arbitrary companion rules}
Adding \(p(s(x))\to s(x)\) destroys termination. For an inert
constant \(c\) (or a fixed object variable), two steps take \(f(s(c))\) to
\(g(s(c),f(s(c)))\); repetition in the new second-child context
gives an infinite ordinary reduction. The absence of companion
\(p\)-rules in the preceding theorem is therefore substantive.
\end{hblemma}

\begin{proofboundary}{def:p12-klop1987-pa-boundaries}{Separate norms, separate equations, separate tree extensions}
The original length annotations, the repaired max/sum equality
annotations and the inherited upper budgets have distinct meanings.
The proofs above establish the exact three PA rules and the expressly
listed seven-rule ordinary presentation, not all ACP rules or all
rewriting modulo commutative/associative choice. Directing choice
commutativity would itself permit an infinite swapping reduction.
The primary CWI appendix's indexed communication, silent-action and
abstraction cases require their own norm and lifting proofs.
Puel's cited report is a different tree-order extension; neither
the PA proof nor the isolated active-origin argument establishes its
general well-quasi-order hypotheses or its unavoidable-set construction.
The distinct 1987 and 1992 presentations retain their own source standing.
\end{proofboundary}

% Source-mapped expansion of the ordered extension in Remark3.11(iii).
\subsection{Global descent of the ordered marker extension}
\label{sec:p12-klop1987-ordered-marker-proof}

Remark~3.11(iii), printed p.22 of the 1987 report, changes both the
objects and the marker operations: sibling order matters, each operator
has a fixed arity, pushing changes just one argument, and simplifying
the left argument is added \parencite{Klop1987CWITutorial}. A proof of
the earlier commutative game does not establish this enlarged game's
well-foundedness. The following supplies that missing argument for
the actual fixed-arity operations. Its dependency is the fully proved
left-to-right LPO of Definitions~\ref{def:p12-ro-data} and
\ref{def:p12-ro-lpo}, with Theorems~\ref{thm:p12-ro-lpo-structural},
\ref{thm:p12-ro-finite-tree} and \ref{thm:p12-ro-lpo-termination}.
The primary Kamin--L\'evy note distinguishes bounded lexicographic
comparison from an unrestricted tuple comparison: the source must also
dominate every target argument
\parencite[handwritten pp.1--8]{KaminLevy1980Attempts}.
Here that domination is established even when different marked copies
have undergone different reductions; it is not assumed from their names.

\begin{hbdefinition}{def:p12-klop1987-om-game}{Fixed-arity ordered trees and all five root operations}
Let \(\Sigma\subseteq\mathbb N\), with one prescribed finite arity
\(a(n)\) for each \(n\in\Sigma\). Plain trees are finite, nonempty
ground terms over this ranked signature. Marked trees use \(n\) or
\(n^*\) at a node of arity \(a(n)\), without permuting children.
For a tree with plain root \(u=n(\vec v)\), let
\(u^*=n^*(\vec v)\). As in source Notation~3.2, this operation is
defined at a plain root; it does not add a second star or silently
interpret a star on an already starred root as idempotent.
Write \(k=a(n)\). The root operations are
\begin{align*}
 \mathrm{PUT}:\quad &n(\vec u)\longrightarrow n^*(\vec u),\\
 \mathrm{SELECT}:\quad &n^*(u_1,\ldots,u_k)\longrightarrow u_i
       &&(1\leq i\leq k),\\
 \mathrm{COPY}:\quad &n^*(\vec u)\longrightarrow
       m(\underbrace{n^*(\vec u),\ldots,n^*(\vec u)}_{a(m)})
       &&(m\in\Sigma,\ m<n),\\
 \mathrm{PUSH}:\quad &n^*(u_1,\ldots,u_i,\ldots,u_k)
       \longrightarrow n(u_1,\ldots,u_i^*,\ldots,u_k),\\
 \mathrm{LEX}:\quad &n^*(u_1,\ldots,u_k)\longrightarrow
       n(u_1^*,\underbrace{n^*(u_1,\ldots,u_k),\ldots,
                      n^*(u_1,\ldots,u_k)}_{k-1})
       &&(k\geq1).
\end{align*}
PUSH requires \(1\leq i\leq k\) and a plain-root \(u_i\);
LEX requires a plain-root \(u_1\), exactly so that the displayed
root-mark operations are defined. Other descendant marks are permitted.
COPY has zero copies exactly when \(a(m)=0\); its copy count is not
free. PUSH neither copies nor deletes its chosen argument. LEX at
arity one has only \(u_1^*\) on its right-hand argument list.

Denote the source's closure under plain-root contexts by
\(\to_{\rm lex}\), and the expressly larger closure under all
finite marked or plain contexts by \(\rightsquigarrow_{\rm lex}\).
The source context clause puts plain labels on the path above the
hole; its fixed sibling subtrees may themselves contain marks.
Every source step is an auxiliary step. The theorem below proves
descent for this larger game, hence for the stated source game.
\end{hbdefinition}

\begin{hblemma}{thm:p12-klop1987-om-bound}{proved=annex:p12-klop1987-om-bound}{One finite ranked signature for an entire run}
If the largest numerical label in an initial marked or plain tree is
\(M\), every descendant under the auxiliary game uses only
\(\Sigma_M=\{n\in\Sigma:n\leq M\}\). This is one fixed finite
ranked signature, with at most \(M+1\) symbols, throughout the run.
The assertion bounds neither the number of nodes nor depth nor the
number of copies or steps. In particular it is not a finite-depth
special case of the ordered extension.
\end{hblemma}

\begin{hbdefinition}{def:p12-klop1987-om-interpretation}{Nonempty weak and strict LPO lower sets}
Fix \(M\) for a run and add a fresh nullary symbol \(\bot\) below
every member of \(\Sigma_M\). It is only a semantic symbol: the
game neither contains nor marks this bottom symbol. Let \(U_M\) be
the ground terms over \(\Sigma_M\cup\{\bot\}\), with total
precedence inherited from the numerical labels and bottom last, and
left-to-right LPO. Write \(>\) for this LPO and \(\geq\) for
equality or \(>\). Define
\[
 D(v)=\{w\in U_M:v\geq w\},\qquad
 D^-(v)=\{w\in U_M:v>w\}.
\]
For a marked game tree define its interpretation recursively by
\begin{align*}
 J(n(u_1,\ldots,u_k))
  &=\bigcup_{b_i\in J(u_i)}D(n(b_1,\ldots,b_k)),\\
 J(n^*(u_1,\ldots,u_k))
  &=\bigcup_{b_i\in J(u_i)}D^-(n(b_1,\ldots,b_k)).
\end{align*}
The product over an empty argument list has one empty tuple.
A lower set means downward closed for \(\geq\). Lower sets here
may be infinite and need not have greatest elements.
\end{hbdefinition}

\begin{hblemma}{thm:p12-klop1987-om-sets}{proved=annex:p12-klop1987-om-sets}{Well-defined interpretations, context monotonicity and finite witnesses}
The interpretation of every finite marked tree is a nonempty lower
set containing \(\bot\). Both constructor interpretations are
monotone under inclusion of each argument set. Consequently replacing
a subtree by one with a smaller interpretation cannot enlarge the
interpretation in any marked or plain context.
Every finite nonempty family of members of one such argument set has
a greatest member belonging to that set, since ground LPO is total.
Thus finitely many argument-choice witnesses can be dominated
coordinatewise within the same argument sets. This uses maxima of
finite families, not a greatest member of an entire lower set.
\end{hblemma}

The added bottom repairs an actual failure of the naive interpretation.
Without it, over a signature with nullary \(c\) and binary \(n>c\),
the least term \(c\) has \(D^-(c)=\varnothing\). The naive
interpretations would then give \(J(n^*(c^*,c))=\varnothing\)
but \(J(c)=\{c\}\), although SELECT of the second child is
\(n^*(c^*,c)\to c\). Thus downward closure alone, allowing empty
argument sets, would not justify operation monotonicity.

\begin{hblemma}{thm:p12-klop1987-om-operations}{proved=annex:p12-klop1987-om-operations}{Every operation decreases the lower-set interpretation}
For each of PUT, SELECT, COPY, PUSH and LEX, and hence for every
\(u\rightsquigarrow_{\rm lex}v\), one has \(J(v)\subseteq J(u)\).
COPY includes its nullary case. The LEX assertion covers every fixed
arity \(k\geq1\), with independent interpretation choices in its
\(k-1\) marked copies. No equality between those choices is required.
\end{hblemma}

\begin{hblemma}{thm:p12-klop1987-om-first-mark}{proved=annex:p12-klop1987-om-first-mark}{The first mark gives strict endpoint descent}
For an entirely plain tree \(t\), \(J(t)=D(t)\). For its root
mark \(t^*\), \(J(t^*)=D^-(t)\). If a first PUT is performed
at any position of an entirely plain tree \(t\), producing \(u\),
then \(J(u)\subseteq D^-(t)\). Therefore every positive finite
auxiliary reduction between plain trees satisfies
\[
 t\rightsquigarrow_{\rm lex}^{+}s\quad\Longrightarrow\quad t>s.
\]
The individual auxiliary steps need only preserve inclusion; the
first mark supplies the strictness retained at the plain endpoint.
\end{hblemma}

\begin{hbtheorem}{thm:p12-klop1987-om-well-founded}{proved=annex:p12-klop1987-om-well-founded}{Global well-foundedness of the ordered fixed-arity marker order}
For every ranked signature \(\Sigma\subseteq\mathbb N\), without
one common global label, arity, size or depth bound, the positive
auxiliary relation restricted to pairs of plain trees is transitive,
irreflexive and well-founded. The same holds for the source relation.
In particular no infinite sequence of plain trees can be connected by
successive positive finite reductions using the ordered LEX extension.
This is the missing global conclusion for Remark~3.11(iii), not just
an orientation of the report's displayed Ackermann pair.
\end{hbtheorem}

This conclusion concerns plain endpoints, not termination of the whole
marked universe. For a binary \(n\), a plain-root \(c\) and any
\(u\), the auxiliary game has the marked cycle
\[
 n^*(c,u)\longrightarrow n(c^*,n^*(c,u))
 \longrightarrow n^*(c^*,n^*(c,u))\longrightarrow n^*(c,u).
\]
The steps are LEX, PUT and SELECT of the second child. Its endpoints
are marked, so it contradicts neither the theorem nor the strict
plain-endpoint interpretation.

\begin{hbtheorem}{thm:p12-klop1987-om-applications}{proved=annex:p12-klop1987-om-applications}{All source Ackermann rules and the ordered associativity example terminate}
The report's three rules
\[
 A(0,x)\to S(x),\qquad A(S(x),0)\to A(x,S(0)),\qquad
 A(S(x),S(y))\to A(x,A(S(x),y))
\]
have positive ordered-marker simulations with fixed arities
\(a(A)=2\), \(a(S)=1\), \(a(0)=0\) and labels
\(A=2>S=1>0\). They terminate on all finite ordinary terms,
over \(\{A,S,0\}\) with variables, not just closed numerical
inputs. The rule
\((x\mathbin{\cdot}y)\mathbin{\cdot}z\to
x\mathbin{\cdot}(y\mathbin{\cdot}z)\) likewise has a positive
ordered-marker simulation with one fixed binary product operator,
and terminates. The annex supplies the actual paths and transports
them through arbitrary rule instances and contexts.
\end{hbtheorem}

\begin{proofboundary}{def:p12-klop1987-om-standing}{What this proof establishes and what remains a different method}
The lower-set argument is an editorial completion of the 1987 marker
presentation, using the separately proved LPO dependency. It is not
claimed to be printed in full in the tutorial, in the undated primary
Kamin--L\'evy scan, or in the distinct 1992 chapter. Its finite-signature
step is justified by the marker run's numerical-label invariant, not
imposed as a restriction on the theorem's infinite ambient signature.
It neither reidentifies ordered fixed-arity trees with the earlier
commutative unranked trees, nor imports a stronger embedding convention
as the report's own definition. Puel's extension for an embedded left
side and the full norm-indexed process-algebra development remain
separate methods and obligations.
\end{proofboundary}


\subsection{Completion: the particular group calculation and algorithm boundary}

\begin{hbconvention}{def:p12-klop1987-group-signature}{The group signature has its own typed roles}
Here \(e\) is a nullary identity symbol, \(I\) is unary inverse,
and \(.\) is binary noncommutative product. This \(I\) is not the
nullary combinator \(I\) of CL. Products retain their parentheses;
associativity is an equation being oriented, not an implicit equation
in first-order matching. The finite signature is the one used in the
source's group tables.
\end{hbconvention}

The source begins with the group equations
\(e.x=x\), \(I(x).x=e\), \((x.y).z=x.(y.z)\).
Its ground word problem asks whether two ground terms are equated by the
equational theory. A sound, equivalent complete ordinary TRS decides it
by comparison of normal forms. CL's equations show why such a terminating
presentation need not exist. The report also distinguishes a decidable word
problem from a finite complete presentation over the \emph{same signature}.
Its string example retains the postfix restriction introduced by the
constant \(o\); its two-sorted multiset example has a single commutation
equation \(\mathsf{ins}(x,\mathsf{ins}(y,X))=
\mathsf{ins}(y,\mathsf{ins}(x,X))\). The natural-number and multiset
zeros there have different sorts, although the printed glyph is shared.

The hand completion first orients the three group equations. The peak
\((I(x).x).z\) gives \(e.z\) and \(I(x).(x.z)\), hence adds
rule 4, \(I(x).(x.z)\to z\). Subsequent source rule numbers are:
\begin{align*}
5:&\ I(I(y)).e\to y,&6:&\ I(e).z\to z,&7:&\ I(I(y)).x\to y.x,\\
8:&\ y.e\to y,&9:&\ I(I(y))\to y,&10:&\ I(e)\to e,\\
11:&\ y.I(y)\to e,&12:&\ y.(I(y).x)\to x,\\
13:&\ x.(y.I(x.y))\to e,&14:&\ y.I(x.y)\to I(x),\\
15:&\ I(x.y)\to I(y).I(x).
\end{align*}
The redundant rules 5 and 7 are removed after 8 and 9; 10 removes 6;
14 removes 13; 15 removes 14. Source Table 12 retains exactly
1, 2, 3, 4, 8, 9, 10, 11, 12, 15. The print of temporary rule 13
has mismatched parentheses; its preceding critical pair determines the
well-formed term displayed above. The report's joining example reduces
\(I(y.I(y))\) to \(I(e)\to e\), or to
\(I(I(y)).I(y)\to y.I(y)\to e\).
Choosing the other direction of rule 15 is not endorsed merely because
it shortens the written expression.

\begin{hbtheorem}{thm:p12-klop1987-group-completion}{proved=annex:p12-klop1987-group-completion}{The actual ten-rule group completion}
The ten rules retained in source Table 12 terminate and are confluent.
Their equational congruence is exactly that generated by source Table
10's three equations. Consequently two finite terms are equal in that
congruence if and only if their final-rule normal forms are identical;
in particular normalization decides the source's ground word problem.
The proof verifies each rule under LPO with precedence
\(I>_\Sigma .>_\Sigma e\) and left-to-right product status,
derives all successive source rule equations from their actual peaks,
and supplies all \(50\) nontrivial critical-pair joining paths together
with the \(15\) literal trivial pairs. These are the \(65\)
unifiable cases among the \(220\) nonvariable overlap candidates.
This concrete use of the existing LPO interface is not substituted for
the source's separate marked-order/Kruskal proof obligation.
\end{hbtheorem}

Theorem 4.2 is the ordinary critical-pair criterion for WCR; Theorem 4.1
adds SN and Newman's lemma for CR. Figure 11 separates disjoint redexes,
redexes inside substituted variables, and overlaps in nonvariable pattern
positions. The variable case must update every repeated occurrence in a
nonlinear left side and every duplicated occurrence on the right; it cannot
be replaced by an unqualified assertion that nested steps commute.
The full proof at Theorem~\ref{thm:p12-critical-pair} treats those exact
ordinary unconditional cases. The concrete group application is verified
in Theorem~\ref{thm:p12-klop1987-group-completion}, including its
nonlinear left sides and every actual nonvariable overlap.

Table 13's six stages are: select an orientable equation or fail; insert its
decreasing rule; normalize existing right sides; add new critical-pair
equations; remove old rules whose left sides contain a new redex;
normalize equations and delete identities. The order of these stages
matters: adding new critical pairs precedes deleting the old reducible
left sides. For an old rule \(l\to r\) with an occurrence of the
new rule's left side, its old/new overlap adds the equation between
\(r\) and the reduct \(l'\) obtained at that occurrence.
Together with the new-rule reduction \(l\to l'\), that equation
preserves the removed rule's equality. Thus an implementation must not
delete before generating the required overlap equation, or discard
that equation with its old rule. The current host's explicit Collapse
transformation packages this dependency by returning a simplified
equation using \emph{retained} rules;
Lemma~\ref{thm:p12-completion-invariant} proves that invariant at its
own stated scope. It is not necessary to invent an omitted equation
step in the original table when its critical-pair stage already supplies it.
An order used on rules with variables must be stable under
substitution as well as contexts, or must orient every rule instance;
context monotonicity alone does not justify termination of all instances.
That additional order standing is explicit, rather than silently
attributed to the table's shorter formulation.

\subsection{Regular systems and parallel redexes}

\begin{hbdefinition}{def:p12-klop1987-regular}{Regularity is linearity plus absence of pattern overlaps}
A rule is left-linear when every variable occurs at most once in its left
side. A TRS is regular in the report's sense when it is left-linear and
non-ambiguous: no two left-hand-side patterns overlap at nonvariable
positions, apart from comparing an occurrence with itself.
This is the ordinary orthogonal regime; the report's term ``regular'' does
not mean a regular language or simply a confluent relation.
The pattern retains the operator occurrences and replaces variable leaves
by distinct holes. Redexes may be nested, but, in a regular system, any
properly nested redex lies wholly within a variable instance of the outer
rule, not within its operator pattern.

Parallel reduction \(t\to_{\parallel}s\) contracts a finite set of
pairwise disjoint redex occurrences. Here the empty set is allowed, giving
reflexivity; this convention explicitly supplies the zero-step arms needed
in joining diagrams. Any such parallel contraction can be sequentialized
in any order, since replacements in disjoint subtrees do not interfere.
\end{hbdefinition}

\begin{hbexample}{def:p12-klop1987-patterns}{The source's overlap and nonlinear failures}
Source Table 14 has
\(F(G(x,S(0)),y,H(z))\to x\),
\(G(x,S(S(0)))\to0\), \(P(G(x,S(0)))\to S(0)\).
The first pattern contains a \(G\)-node, but cannot overlap the second
rule there: \(S(0)\) and \(S(S(0))\) do not unify. The other heads
are distinct. Figures 12--13 identify those non-overlapping operator
patterns, including variable-instance boundaries. Figure 14 contrasts
the self-overlap of \(L(L(x))\to0\), the root overlap of
\(F(0,x,y)\to0\) with \(F(x,1,y)\to1\), and the two root
rules \(\mathsf{or}(x,y)\to x,\mathsf{or}(x,y)\to y\).

Theorem~\ref{thm:p12-critical-pair} gives WCR, but not CR, for the
nonlinear source example
\(D(x,x)\to E,C(x)\to D(x,C(x)),A\to C(A)\).
Indeed,
\(C(A)\to D(A,C(A))\to D(C(A),C(A))\to E\), whereas
\(C(A)\twoheadrightarrow C(E)\). From \(C(E)\), every term has
one remaining \(C(E)\) under zero or more \(D(E,\ )\)-contexts;
the repeated-variable test never succeeds. Thus it cannot reach \(E\).
Absence of critical pairs does not remove this nonlinear obstruction.
\end{hbexample}

\begin{hbtheorem}{thm:p12-klop1987-regular-confluence}{proved=annex:p12-klop1987-regular-confluence}{The report's regular confluence theorem}
For a regular TRS, disjoint parallel reduction has the diamond property:
if \(t\to_{\parallel}u\) and \(t\to_{\parallel}v\), then some
\(w\) satisfies \(u\to_{\parallel}w\) and
\(v\to_{\parallel}w\). Consequently the ordinary rewrite relation
is CR, without an SN hypothesis. This proves the regular case of
Theorem 5.3 by the residual argument underlying Figures 12--15.
The report attributes the general regular confluence route to Rosen (1973).
\end{hbtheorem}

Theorem 5.5 states Huet's broader condition: in a left-linear TRS, when
each oriented critical pair is connected by a disjoint parallel reduction
in the stated direction, parallel reduction is strongly confluent and
ordinary reduction is CR. This condition is stronger than mere joinability
of critical pairs; it is not Theorem 4.2 with SN silently removed.

\begin{hbexample}{def:p12-klop1987-regular-examples}{CL, recursive program schemes, and nonregular extensions}
CL's three left sides are left-linear and have non-overlapping operator
patterns. Its confluence follows from
Theorem~\ref{thm:p12-klop1987-regular-confluence} despite its infinite
reductions. A recursive program scheme has disjoint finite operator sets
\(\mathcal F=\{F_1,\ldots,F_n\}\) and
\(\mathcal G=\{G_1,\ldots,G_k\}\), each with fixed arities, and
exactly one rule \(F_i(x_1,\ldots,x_{m_i})\to t_i\) for each
\(F_i\), with distinct displayed variables. Its right side uses only
those variables and operators of \(\mathcal F\cup\mathcal G\).
These root patterns cannot overlap, so every such scheme is regular.
The source example is
\[
 F_1(x)\to G_1(x,F_1(x),F_2(x,x)),\qquad
 F_2(x,y)\to G_2(F_2(x,x),F_1(G_3)).
\]
This remains regular despite duplication and erasure on the right.

Both equality-test extensions of Tables 6--7 are nonregular: both
\(D(x,x)\) and \(Dxx\) repeat \(x\). The binary extension's
confluence uses disjoint-sum modularity, not regular confluence;
the exact factor and sum proofs are
Theorem~\ref{thm:p12-klop1987-ds-applications}.
The applicative extension is reported nonconfluent on p.16.
SKIM's numeral instances are left-linear ground rules, rather than
repeated-variable patterns.
\end{hbexample}

\begin{hbdefinition}{def:p12-klop1987-ne-win}{Non-erasure and innermost normalization}
NE means \(\operatorname{Var}(l)=\operatorname{Var}(r)\) in every rule.
It does not forbid duplication. An innermost redex has no proper redex
inside it. A term is WIN when it has a normal form reachable by a finite
innermost reduction; a system is WIN when every term is.
The report states both system-wide and term-local forms of its results.
\end{hbdefinition}

\begin{hbtheorem}{thm:p12-klop1987-ne-normalization}{proved=annex:p12-klop1987-ne-normalization}{The report's non-erasing normalization equivalence}
In a regular NE system, a term is WN if and only if it is SN
(Theorem 5.8, p.32). The equivalence also holds system-wide.
It does not say that arbitrary reductions of a WN term terminate
in an erasing regular system.
\end{hbtheorem}

\begin{hbtheorem}{thm:p12-klop1987-win-normalization}{proved=annex:p12-klop1987-win-normalization}{O'Donnell's innermost equivalence as reported by Klop}
In a regular system, a term is WIN if and only if it is SN; the
equivalence also holds system-wide. This is Theorem 5.9, p.32,
attributed to O'Donnell (1977). Its innermost premise is different from
non-erasure and from mere existence of a normal form.
\end{hbtheorem}

For the unary encodings of a non-overlapping semi-Thue system, NE is
automatic. The report also gives a direct explanation: their strict
one-step diamond property \(\mathrm{WCR}^{1}\), unlike
\(\mathrm{WCR}^{\leq1}\), joins both arms by genuine steps;
together with WN it forces SN. This is Exercise 1.7(15), not a claim that
all subcommutative WN systems terminate. The self-loop rule \(a\to a\)
adjoined to \(a\to b\), with \(b\) normal, is a counterexample
to that weaker assertion.

\subsection{Strategies: normalization does not imply cofinality}

\begin{hbdefinition}{def:p12-klop1987-strategy}{Computable choices and their two success properties}
A one-step strategy is a computable map
\(\mathbb F:\operatorname{Ter}(R)\to\operatorname{Ter}(R)\)
fixing normal forms and satisfying \(t\to\mathbb F(t)\) otherwise.
A many-step strategy instead requires \(t\to^+\mathbb F(t)\).
The iteration is \(t,\mathbb F(t),\mathbb F^2(t),\ldots\).
It is normalizing if every term with a normal form reaches one in this
iteration. It is cofinal if, for every descendant \(s\) of \(t\),
some iterate is reachable from \(s\).
The latter property addresses divergent terms as well: it does not assert
that an infinite sequence has a finite normal form.
\end{hbdefinition}

\begin{hbtheorem}{thm:p12-klop1987-cofinal-normalizing}{proved=annex:p12-klop1987-cofinal-normalizing}{Cofinality forces the right finite result}
A cofinal strategy is normalizing. In a confluent system, whenever a
strategy reaches a normal form it is the unique normal form of its input.
Cofinality is strictly more demanding than normalizing behavior on inputs
which happen to normalize.
\end{hbtheorem}

\begin{hbdefinition}{def:p12-klop1987-five-strategies}{The five strategies of Definition 6.3}
Leftmost-innermost chooses the leftmost minimal redex in the term tree;
parallel-innermost contracts every minimal redex. Leftmost-outermost
chooses the leftmost maximal redex; parallel-outermost contracts every
maximal redex. Minimal and maximal here refer to containment of redex
occurrences, not the lexicographic position order alone. A leftmost
choice uses the tree's left-to-right argument order.

Full substitution, also called Kleene or Gross--Knuth reduction in the
report, develops all redexes present in the input, including nested ones.
It does not indefinitely contract every newly created redex. In the
regular regime a selected inner redex lies in an outer variable instance;
after outer contraction it has one residual at each occurrence of that
variable on the right, or none if erased. All these residuals must be
developed. This specifies ``simultaneously'' without requiring that
the original nested occurrences be disjoint.
\end{hbdefinition}

\begin{hblemma}{thm:p12-klop1987-developments}{proved=annex:p12-klop1987-developments}{Finite developments specify full substitution}
In a regular TRS, every development of a fixed finite set of redex
occurrences terminates, and every complete development has the same result.
Taking that set to be all redexes of the finite input gives a well-defined
many-step operation fixing normal forms. A computable presentation of the
rules and matching supplies the effective strategy; computability is not
automatic for an arbitrary set-theoretic rule family.
\end{hblemma}

\begin{hbdefinition}{def:p12-klop1987-fair-traces}{Residual fairness, not merely an infinite choice sequence}
A residual trace through a reduction chooses one surviving residual
occurrence at every successive term, starting at any finite stage.
A fair reduction has no infinite redex-residual trace. Erasure or actual
contraction ends a trace; duplication permits finitely many choices for
its next member. A newly created redex begins a new trace and is not
identified with the contracted redex merely because the printed terms
are equal. Full substitution is sequentialized by completing one finite
development round before starting the next.
A complete fair run is either an infinite fair reduction or a finite
reduction ending at a normal form; finiteness alone does not complete a run.
\end{hbdefinition}

\begin{hbtheorem}{thm:p12-klop1987-fair-cofinal}{proved=annex:p12-klop1987-fair-cofinal}{Fair reductions and full-substitution cofinality}
In a regular finite-term TRS, a complete fair run is cofinal in its input's
reduction graph. In particular the full-substitution strategy is cofinal
and hence normalizing, without NE or SN assumptions. The assertion for
a finite reduction includes the requirement that its endpoint be normal;
an arbitrarily truncated finite prefix is not a cofinal strategy run.
This supplies the double-plus row of the source's strategy table.
\end{hbtheorem}

\begin{hbdefinition}{def:p12-klop1987-outermost-fair}{Outermost traces and finite projections}
A reduction is outermost-fair if it has no infinite residual trace whose
members remain outermost at every stage. It can have infinite traces
that repeatedly cease to be outermost; those are not excluded by this
definition. Projection through a chosen step tracks that step's disjoint
residual set along the other reduction, develops that set at each finite
stage, and connects the endpoints by the regular commuting position cases.
Empty projected segments are recorded as empty, not as extra rewrite steps.
\end{hbdefinition}

\begin{hblemma}{thm:p12-klop1987-outermost-projection}{proved=annex:p12-klop1987-outermost-projection}{The two ordinary outermost projection properties}
In a regular finite-term TRS, projection through a fixed step preserves
outermost fairness. If the original reduction has infinitely many
outermost contractions, its projection is infinite. Every infinite
outermost-fair reduction has infinitely many outermost contractions.
These are the ordinary-TRS mechanisms in the tutorial-cited
Bergstra--Klop appendix, specialized without conditions or bound variables
\parencite[pp.352--360]{BergstraKlop1986StrategyAppendix}.
\end{hblemma}

\begin{hbdefinition}{def:p12-klop1987-left-normal}{The additional left-normal restriction}
A regular TRS is left-normal when every constant or function-symbol
occurrence in each left side precedes every variable occurrence in its
linear prefix notation. CL, when restored to binary application notation,
and the recursive program schemes above have this property.
The pattern \(F(x,B)\) does not: the constant \(B\) follows \(x\).
Left-normality is an additional restriction, not another name for regularity.
\end{hbdefinition}

\begin{hbtheorem}{thm:p12-klop1987-strategy-results}{proved=annex:p12-klop1987-strategy-results}{The source's strategy table, pp.34--35}
On regular TRSs, parallel-outermost reduction is normalizing. On left-normal regular TRSs,
leftmost-outermost reduction is also normalizing. Innermost strategies
are not in general normalizing even in the left-normal regular regime.
Neither outermost strategy is in general cofinal, including in that
restricted regime. The report refers the two outermost normalization
proofs to O'Donnell (1977), and the parallel-outermost proof also to the
appendix of Bergstra--Klop (1986).
\end{hbtheorem}

The source uses minus for a failure of normalization in general, plus
for normalizing but not necessarily cofinal, and double-plus for cofinal.
Thus its table records, in order, two minus rows for the innermost
strategies, minus/plus for leftmost-outermost, plus/plus for
parallel-outermost, and double-plus/double-plus for full substitution.
It is not a termination table for all terms.

\begin{hblemma}{thm:p12-klop1987-leftmost-needed}{proved=annex:p12-klop1987-leftmost-needed}{A needed-redex dependency in the effective left-normal regime}
In a regular left-normal TRS, the leftmost redex of a normalizable
non-normal term is needed: every reduction to a normal form contracts
its residual. The occurrence can be found by finite position search
under the same effective-matching assumption as the strategies above.
This is a special regime, not a general computable selector for all
regular systems or a decision test for normalizability.
\end{hblemma}

\begin{hbexample}{def:p12-klop1987-strategy-failures}{The source's leftmost obstruction and the erasure obstruction}
Huet--L\'evy's example on p.35 is the regular system
\(F(x,B)\to D,A\to B,C\to C\).
The term \(F(C,A)\to F(C,B)\to D\) has a normal form.
But leftmost-outermost repeatedly contracts its \(C\)-occurrence,
leaving \(F(C,A)\) unchanged. Parallel-outermost instead contracts
both \(C\) and \(A\), then the root, and reaches \(D\).
There is no contradiction with the left-normal theorem because the
\(F\)-pattern is not left-normal.

An explicit erasure example, separate from the source's displayed
counterexample, is \(Kxy\to x,a\to a\) and \(Kba\), where
\(b\) is normal. This regular left-normal term reaches \(b\)
by an outermost step, but either innermost strategy loops at \(a\).
It illustrates why Theorem~\ref{thm:p12-klop1987-ne-normalization}
retains NE and Theorem~\ref{thm:p12-klop1987-win-normalization}
requires WIN rather than assuming WN alone.
\end{hbexample}

\begin{hbtheorem}{thm:p12-klop1987-omega}{proved=annex:p12-klop1987-omega}{The source's CL failure of outermost cofinality}
Put \(I_t=SKt\) and \(\Omega_t=SI_tI_t(SI_tI_t)\) in CL.
Then \(\Omega_{II}\twoheadrightarrow\Omega_I\), but
\(\Omega_I\not\twoheadrightarrow\Omega_{II}\).
The parallel-outermost orbit of \(\Omega_{II}\) is periodic and is
not cofinal. This is the source's p.35 example; normalizing and cofinal
are therefore genuinely different strategy requirements even for CL,
which is regular and left-normal.
\end{hbtheorem}

The source's stronger outermost failure uses CL. Put
\(I_t=SKt\), so \(I_t x\twoheadrightarrow x\), and
\(\Omega_t=SI_tI_t(SI_tI_t)\).
Writing \(W_t=SI_tI_t\), its parallel-outermost cycle is
\[
 W_tW_t\to I_tW_t(I_tW_t)
 \to_{\parallel}KW_t(tW_t)(KW_t(tW_t))
 \to_{\parallel}W_tW_t.
\]
With \(t=II\), the source notes
\(\Omega_{II}\twoheadrightarrow\Omega_I\), while
\(\Omega_I\not\twoheadrightarrow\Omega_{II}\); consequently the
cyclic parallel-outermost orbit is not cofinal. The nonreachability claim
has a separate invariant proof in Annex~\ref{annex:p12-klop1987-omega};
it is not proved merely by displaying the cycle. It supplies the cofinality distinction rather than another
finite-normal-form failure.

The final source paragraph attributes to Huet--L\'evy the theory of
needed redexes and strongly sequential regular TRSs. A needed occurrence
has a residual contracted in every reduction to a normal form. For a
normalizable non-normal term the report asserts existence; its wording
must not be applied to a normal form, which has no redex at all, or made
informative by vacuity on a term with no normal form. It states that a
general computable needed-redex selector is unavailable, whereas strong
sequentiality is decidable and supplies an efficient selector and hence
a normalizing strategy. The tutorial stops before developing that theory;
no decision procedure for strong sequentiality is specified here.

\subsection{The exact source/version boundary}

The title leaf of this report announces intended publication in
\emph{Bulletin of the European Association for Theoretical Computer
Science}, no.32, 1987. That statement is not an announcement of the
Oxford handbook chapter published in 1992. The report's own abstract
excludes conditional and equational TRSs. The publisher chapter's
separately retained abstract includes E-unification, strongly sequential
orthogonal TRSs, and conditional TRSs. Shared author, subject, and theorem
names therefore justify a relationship for comparison, not replacement
of one presentation by the other. A full page-by-page substantive
variant crosswalk with the publisher chapter is not supplied by this
report or by this reconstruction. In particular, the historical one-rule
termination question retains its report date; later chapter statements
have their own source/version attribution.

\begin{proofboundary}{def:p12-klop1987-proof-boundary}{The mathematical boundaries of the residual arguments}
The regular diamond and strategy arguments concern ordinary first-order
regular systems on finite terms. They do not prove Huet's broader
left-linear parallel-critical-pair criterion merely from ordinary
joinability, and they do not extend to conditional rules, rewriting
modulo equations, infinite terms, or bound-variable systems without
their separate residual hypotheses. Non-erasure and innermostness are
different sufficient premises for local WN/SN and WIN/SN equivalences;
neither may be omitted in favor of WN alone.
The needed-redex selector proved here uses left-normality. The source's
general needed-redex and strongly sequential assertions retain their
separate cited standing, not that selector's proof or effective method.
The unranked commutative marked-tree well-foundedness route instead
depends on its embedding and Kruskal arguments; a finite-development
proof is not a substitute for them. The ordered fixed-arity LEX extension
has its separate nonempty-LPO-lower-set proof and proves descent only
between plain endpoints, not termination of its marked auxiliary universe.
Similarly, the group completion requires its actual decreasing
order and every critical pair, rather than only the abstract completion
criterion or a differently presented host example.
\end{proofboundary}
